\documentclass[11pt,a4paper]{article}
\input{preamble}
\usepackage{graphicx}
\usepackage{xcolor}
\graphicspath{{figures/dynamical-systems/}}
\title{Dynamical Systems}
\begin{document}
\maketitle
\thispagestyle{empty}

\noindent
\underline{Dynamical systems} -- mathematical entities described by a set of
equations for a set of variables w.r.t. a time-like variable.

\medskip
\noindent
\underline{State space} or \underline{phase space}.

\medskip
\noindent
$\ast$ \underline{Continuous time} $\Rightarrow$ set of ODEs

\medskip
\noindent
Lotka--Volterra for rabbits $(r)$ \& sheep $(s)$
\[
\dot r=r(a-br-cs),\qquad \dot s=s(d-er-fs),\qquad a,b,\ldots,f>0.
\]
\[
\text{State space}\qquad [0,\infty)\times[0,\infty).
\]

\noindent
e.g.\ \underline{Logistic map}\qquad
$x_{n+1}=\mu x_n(1-x_n),\quad 0\leq\mu\leq1$

\medskip
\noindent\underline{Questions to address}
\begin{itemize}
    \setlength{\itemsep}{0pt}
    \item[$\ast$] Long-term behavior of the system?
    \item[$\ast$] Are there periodic solutions
    \item[$\ast$] Stability of solutions, chaos
    \item[$\ast$] Effect of system parameter variation
\end{itemize}

\noindent\fbox{Example}\quad 1D system\qquad $\dot x=x(1-x^2)$
\[
\text{Solution}\qquad x(t)=\pm\left(1+2e^{2(t_0-t)}\right)^{-1/2}
\]
\noindent
\begin{minipage}[c]{0.66\textwidth}
    \includegraphics[width=\linewidth]{one-dimensional.png}
\end{minipage}\hfill
\begin{minipage}[c]{0.32\textwidth}
    \centering
    \underline{as $t\to\infty$}
    \[
    \begin{aligned}
    x(0)>0&\Rightarrow x(t)\to1\\
    x(0)<0&\Rightarrow x(t)\to-1\\
    x(0)=0&\Rightarrow x(t)=0
    \end{aligned}
    \]
\end{minipage}

\newpage
\noindent\fbox{Example}\quad 2D, Lotka--Volterra

\medskip
\noindent
\begin{minipage}[c]{0.43\textwidth}
    \[
    \begin{aligned}
    \dot r&=r(3-r-s)\\
    \dot s&=s(2-r-s)
    \end{aligned}
    \]
    Can get $r(s)\Rightarrow t(s)$

    \medskip
    $\ast$ \underline{Find fixed points}
    \[
    \dot r,\dot s=0:\quad (0,0),(3,0),(0,2)
    \]
    $\ast$ \underline{Look at $\dot r,\dot s$ in state space}
\end{minipage}\hfill
\begin{minipage}[c]{0.55\textwidth}
    \includegraphics[width=\linewidth]{lotka-volterra.png}
\end{minipage}

\[
\text{As }t\to\infty\qquad
(r,s)\to
\begin{cases}
    (3,0),&r>0,\\
    (0,2),&r=0,\ s>0,\\
    (0,0),&r=s=0.
\end{cases}
\]

\bigskip
\noindent\fbox{Example}\quad \underline{predator-prey}

\medskip
\noindent
\begin{minipage}[c]{0.43\textwidth}
    sheep $(s)$ \& wolves $(w)$
    \[
    \begin{aligned}
    \dot w&=w(-a+bs)\\
    \dot s&=s(c-dw)
    \end{aligned}
    \]
    \[
    \underline{a,\ldots,d>0}
    \]
    Fixed points
    \[
    (0,0),\quad (a/b,c/d)
    \]
\end{minipage}\hfill
\begin{minipage}[c]{0.55\textwidth}
    \includegraphics[width=\linewidth]{predator-prey.png}
\end{minipage}

\[
(s,w)\to
\begin{cases}
    \text{periodic around }(a/b,c/d),&s>0,\ w>0,\\
    (0,0),&s=0,\ w\geq0,\\
    (\infty,0),&s>0,\ w=0
\end{cases}
\qquad\text{as }t\to\infty
\]

\newpage
\noindent\fbox{Example}
\[
\begin{aligned}
\dot x&=-y+\varepsilon x(\mu-x^2-y^2)\\
\dot y&=\phantom{-}x+\varepsilon y(\mu-x^2-y^2)
\end{aligned}
\]

\noindent
Switch to polars:
\[
r=\sqrt{x^2+y^2},\qquad \theta=\operatorname{atan}\frac{y}{x},
\]
\[
\dot r=\frac{x\dot x+y\dot y}{r},\qquad
\dot\theta=\frac{x\dot y-y\dot x}{r^2}
\quad\Rightarrow
\]
\[
\underline{\dot r=\varepsilon r(\mu-r^2)},\qquad
\underline{\dot\theta=1}
\]

\begin{center}
    \includegraphics[width=\textwidth]{polar-cases.png}
\end{center}

\begin{center}
    \textbf{End of lecture 1}
\end{center}


\newpage
\noindent
\underline{Poincare} map in 2D for a 3D dynamical system:

\begin{center}
    \includegraphics[width=\textwidth]{lecture2-poincare-maps.png}
\end{center}

\noindent
Can also consider maps as discretized ODEs.

\bigskip
\noindent\underline{Plan for the course}
\[
\begin{array}{l@{\quad}l}
\left.\begin{array}{rll}
(3)&\S 1:&\text{\makebox[36mm][l]{Basic definitions}}\\[6pt]
(3)&\S 2:&\text{Fixed points}\\[6pt]
(2)&\S 3:&\text{Stability}\\[6pt]
(5)&\S 4:&\text{Periodic orbits}\\[6pt]
(5)&\S 5:&\text{Bifurcations}
\end{array}\right\}&\text{ODEs}\\[8pt]
\left.\begin{array}{rll}
(3)&\S 6:&\text{\makebox[36mm][l]{Maps}}\\[6pt]
(3)&\S 7:&\text{Chaos}
\end{array}\right\}&\text{Maps}
\end{array}
\]

\newpage
\noindent\S 1.\ \underline{Definitions}

\bigskip
\noindent
\begingroup\setlength{\unitlength}{1pt}\begin{picture}(24,18)(0,-4)\put(12,0){\oval(24,18)}\put(12,0){\makebox(0,0){1.1}}\end{picture}\endgroup\quad Define a \underline{State Space} (\underline{Phase Space}) $E\subseteq\mathbb{R}^n$.
The state of the system is described by $\bar x\in E$.
The state evolves in time according to an ODE of the form
\[
\underline{\dot{\bar x}=\frac{d\bar x}{dt}=\bar f(\bar x,t)},
\qquad\text{where }\bar f:E\times\mathbb{R}\to E\text{ is a vector field.}
\]
\noindent
If $\dfrac{\partial\bar f}{\partial t}=0$ the system is \underline{autonomous}.

\medskip
\noindent
A non-autonomous system reduces to an autonomous by defining
$\bar y=(\bar x,t)\in E\times\mathbb{R}$ so that
\[
\underline{\dot{\bar y}=\begin{pmatrix}\bar f(\bar y)\\1\end{pmatrix}}.
\]
\noindent
Thus, can study only autonomous systems.

\medskip
\noindent
A \underline{$n^{\text{th}}$ order ODE}
\[
\frac{d^n x}{dt^n}=g\left(x,\frac{dx}{dt},\ldots,\frac{d^{n-1}x}{dt^{n-1}}\right)
\]
can be reduced to a standard form with $n$ equations by introducing
\[
\bar y=\left(x,\frac{dx}{dt},\ldots,\frac{d^{n-1}x}{dt^{n-1}}\right)
=(y_1,y_2,\ldots,y_n)
\]
so that
\[
\dot{\bar y}=\begin{pmatrix}
y_2\\y_3\\\vdots\\y_n\\g(y_1,y_2,\ldots,y_n)
\end{pmatrix}.
\]
\noindent
Similarly to autonomous systems, will only consider maps of the form
\[
\underline{\bar x_{n+1}=\bar F(\bar x_n)}
\qquad\text{and not}\qquad\bar F(\bar x_n,n).
\]

\newpage
\noindent\begingroup\setlength{\unitlength}{1pt}\begin{picture}(24,18)(0,-4)\put(12,0){\oval(24,18)}\put(12,0){\makebox(0,0){1.2}}\end{picture}\endgroup\quad\underline{Initial value problem (IVP)}

\bigskip
\noindent
We would like to explore an \underline{initial value problem} for
$\underline{\dot{\bar x}=\bar f(\bar x,t)}$:

\medskip
\noindent
\begin{tabular}{@{\vrule width 0.5pt\hspace{8pt}}p{\dimexpr\textwidth-10pt\relax}@{}}
Given the initial value of $\bar x$: $\underline{\bar x(t_0)=\bar x_0}$,
$\bar x_0\in E$, $t\in I\subseteq\mathbb{R}$, find a differentiable function
$\bar x(t)$ for $t\in I$ which remains in $E$ for $t\in I$ and satisfies the
initial condition and the ODE.
\end{tabular}

\bigskip
\noindent
For an autonomous system can define solution in terms of a \underline{flow $\phi_t$}.

\medskip
\noindent\fbox{Definition}\quad (\underline{Flow}) $\phi_t(\bar x)$ such that
$\phi_t(\bar x_0)$ is the solution of $\dot{\bar x}=\bar f(\bar x)$ starting at
$\bar x_0$ when $t=0$ is called the flow through $\bar x_0$ at $t=0$.

\medskip
\noindent Thus,
\[
\phi_0(\bar x_0)=\bar x_0,
\qquad
\phi_s\bigl(\phi_t(\bar x_0)\bigr)=\phi_{s+t}(\bar x_0).
\]
\noindent
Does this solution \underline{exist}? Is it \underline{unique}?

\bigskip
\noindent\fbox{Theorem}\quad (Cauchy--Peano)

\medskip
\noindent
If $\bar f(\bar x,t)$ is continuous and $|\bar f|<M$ in the domain
\[
D:\bigl\{|t-t_0|<\alpha,\;|\bar x-\bar x_0|<\beta\bigr\}
\]
then the IVP has a solution for
$|t-t_0|<\underline{\min(\alpha,\beta/M)}$.

\bigskip
\noindent
But \underline{uniqueness} requires stronger conditions on $f$.

\newpage
\noindent
\begin{minipage}[t]{0.60\textwidth}
\vspace{0pt}
\fbox{Example}\quad \underline{Unique solution}
\[
\dot x=|x|,\qquad x(t_0)=x_0
\]
Solution
\[
x(t)=\begin{cases}
x_0e^{t_0-t},&x_0<0,\\
0,&x_0=0,\\
x_0e^{t-t_0},&x_0>0.
\end{cases}
\]
Note that $f$ is \underline{continuous} but \underline{not differentiable}.
\end{minipage}\hfill
\begin{minipage}[t]{0.37\textwidth}
\vspace{0pt}
\includegraphics[width=\linewidth]{lecture2-unique-solution.png}
\end{minipage}
\par\noindent\hfill$\blacktriangleleft$

\medskip
\noindent
\begin{minipage}[t]{0.60\textwidth}
\vspace{0pt}
\fbox{Example}\quad \underline{Non-unique solution}
\[
\dot x=\begin{cases}
\sqrt{x},&x\geq0,\\
0,&x<0,
\end{cases}
\qquad\underline{x(0)=0}
\]
Can check that
\end{minipage}\hfill
\begin{minipage}[t]{0.37\textwidth}
\vspace{0pt}
\includegraphics[width=\linewidth]{lecture2-nonunique-vector-field.png}
\end{minipage}
\[
x(t)=\begin{cases}
0,&t<\tau,\\
\frac14(t-\tau)^2,&t\geq\tau
\end{cases}
\qquad\text{for any }\tau\geq0
\]
\noindent
\begin{minipage}[t]{0.60\textwidth}
\vspace{0pt}
satisfies the equation and initial condition. But $\tau$ can be anything so a
\underline{family} of solutions!

\medskip
Here again $f$ is \underline{continuous} and \underline{not differentiable}.
\end{minipage}\hfill
\begin{minipage}[t]{0.37\textwidth}
\vspace{0pt}
\includegraphics[width=\linewidth]{lecture2-nonunique-solutions.png}
\end{minipage}
\par\noindent\hfill$\blacktriangleleft$

\medskip
\noindent
Why is the difference? Because $\dfrac{d}{dx}(\sqrt{x})$ is not bounded at $0$!

\medskip
\noindent
To guarantee uniqueness, need $f$ to be not just continuous, but also
\underline{Lipschitz}.

\newpage
\noindent\fbox{Definition}\quad \underline{Lipschitz property}

\medskip
\noindent
A function $\bar f$ defined on a subset of $\mathbb{R}^n$ satisfies Lipschitz
condition at a point $\bar x_0$ with Lipschitz constant $L$ if $\exists\,(L,a)$
such that $\forall\,\bar x,\bar y$: $|\bar x-\bar x_0|<a$,
$|\bar y-\bar x_0|<a$ we have
\[
|\bar f(\bar x)-\bar f(\bar y)|\leq L|\bar x-\bar y|.
\]

\medskip
\noindent
Note that Differentiable $\to$ Lipschitz $\to$ Continuous

\bigskip
\noindent\fbox{Theorem}\quad
\underline{Uniqueness/Picard--Lindel\"of/Cauchy--Lipschitz}

\medskip
\noindent
Consider an IVP $\dot{\bar x}=\bar f(\bar x,t)$ with $\bar x=\bar x_0$ at $t=t_0$.
If $\bar f$ satisfies a Lipschitz condition at $\bar x_0$ then the solution
$\phi_{t-t_0}(\bar x_0)$ \underline{exists} and is \underline{unique} and
\underline{continuous} in a neighborhood of $(\bar x_0,t_0)$.

\bigskip
\noindent
However, in general uniqueness \& continuity do not necessarily mean existence
of solution for all time!

\bigskip
\noindent\fbox{Example}\quad \underline{Finite time blowup}
\[
\dot x=x^2,\qquad x(0)=1,\qquad\text{solution is }x(t)=\frac{1}{1-t},
\]
\noindent
blows up to $\infty$ at $t\to1$.

\bigskip
\noindent
In what follows, unless stated otherwise, we will consider only
\underline{differentiable} functions $f$.

\newpage
\noindent\S 1.3.\ \underline{Orbits, invariant sets and limit sets}

\bigskip
\noindent
Consider a flow $\phi_t(\bar x_0)$.

\medskip
\noindent
\begin{minipage}[c]{0.58\textwidth}
\fbox{Definition}

\medskip
$\ast$ The \underline{orbit/trajectory} of $\phi_t$ through $\bar x_0$ is the
set
\[
O(\bar x_0)=\{\phi_t(\bar x_0):-\infty<t<\infty\}.
\]
$\ast$ The \underline{forward orbit} of $\phi_t$ through $\bar x_0$ is the
set
\[
O^+(\bar x_0)=\{\phi_t(\bar x_0):0\leq t\};
\]
analogously, \underline{backward orbit} $O^-(\bar x_0)$ if for $t\leq0$.
\end{minipage}\hfill
\begin{minipage}[c]{0.39\textwidth}
\includegraphics[width=\linewidth]{lecture3-orbit.png}
\end{minipage}

\bigskip
\noindent\fbox{Definition}

\medskip
\noindent
A set of points $\mathcal M\subseteq E$ is \underline{invariant} under $f$ if
\[
\bar x\in\mathcal M\quad\Rightarrow\quad O(\bar x)\subseteq\mathcal M.
\]
\noindent
Clearly, any orbit is invariant.

\bigskip
\noindent\fbox{Definition}\quad\underline{Fixed point}

\medskip
\noindent
The point $\bar x_0$ is a \underline{fixed point} (equilibrium, stationary,
critical point) if $f(\bar x_0)=0$. Then $\bar x=\bar x_0$ for any $t$ and
$O(\bar x_0)=\bar x_0$.

\bigskip
\noindent\fbox{Definition}\quad\underline{Periodic point \& orbit}

\medskip
\noindent
A point $\bar x_0$ is a \underline{periodic point} with period $T$ if
$\phi_T(\bar x_0)=\bar x_0$ for some $T>0$ but
$\phi_t(\bar x_0)\neq\bar x_0$ for $0<t<T$.

\medskip
\noindent
The set $\{\phi_t(\bar x_0):0\leq t<T\}$ is the \underline{periodic orbit}
through $\bar x_0$.

\newpage
\noindent\fbox{Example}\quad\underline{Periodic orbit}
\[
\begin{pmatrix}\dot x\\[2pt]\dot y\end{pmatrix}
=
\begin{pmatrix}-y^3\\[2pt]x^3\end{pmatrix}
\quad\Rightarrow\quad
\underline{x^4+y^4=\operatorname{cst}(t)}
\]

\begin{center}
\includegraphics[width=0.42\textwidth]{lecture3-periodic-orbits.png}

not a limit cycle!
\end{center}

\noindent\fbox{Definition}\quad\underline{Limit cycle}

\medskip
\noindent
A \underline{limit cycle} is an \underline{isolated} periodic orbit, i.e.
there are no other periodic orbits in its sufficiently small neighborhood.

\bigskip
\noindent\fbox{Example}\quad\underline{Limit cycle}

\medskip
\noindent
\begin{minipage}[c]{0.55\textwidth}
\[
\begin{pmatrix}\dot x\\[2pt]\dot y\end{pmatrix}
=
\begin{pmatrix}
-y+x(1-x^2-y^2)\\[2pt]
x+y(1-x^2-y^2)
\end{pmatrix}
\quad\Rightarrow
\]
\[
\underline{\dot r=r(1-r^2)},
\qquad
\underline{\dot\theta=1}.
\]
A fixed point $(0,0)$ and a \underline{limit cycle} $r=1$.
\end{minipage}\hfill
\begin{minipage}[c]{0.42\textwidth}
\includegraphics[width=\linewidth]{lecture3-limit-cycle.png}
\end{minipage}

\bigskip
\noindent\fbox{Definition}\quad\underline{Homoclinic orbit}

\medskip
\noindent
If $\bar x_0$ is a fixed point and $\exists\,\bar y\neq\bar x_0$ such that
$\phi_t(\bar y)\to\bar x_0$ as $t\to\infty$ and $t\to-\infty$, then
$O(\bar y)$ is a \underline{homoclinic orbit}.

\begin{center}
\includegraphics[width=0.78\textwidth]{lecture3-homoclinic-orbit.png}
\end{center}

\newpage
\noindent\fbox{Definition}\quad\underline{Heteroclinic orbit}

\medskip
\noindent
If $\bar x_0$ and $\bar x_1\neq\bar x_0$ are fixed points and
$\exists\,\bar y$ such that $\phi_t(\bar y)\to\bar x_0$ as $t\to\infty$ and
$\phi_t(\bar y)\to\bar x_1$ at $t\to-\infty$, then $O(\bar y)$ is a
\underline{heteroclinic orbit}.

\begin{center}
Note that it takes $\infty$ time to leave $\bar x_1$ and reach $\bar x_0$!

\includegraphics[width=0.82\textwidth]{lecture3-heteroclinic-orbit.png}
\end{center}

\noindent
In higher dimensions more complex invariant sets are possible.

\bigskip
\noindent\fbox{Example}\quad\underline{2-torus in 3D}

\medskip
\noindent
$\theta_1,\theta_2$ coordinates on the surface of a torus in 3D,
$\dot\theta_1=\omega_1$, $\dot\theta_2=\omega_2$ and $\omega_1,\omega_2$ are
not in a ratio, $\Rightarrow$ trajectory will fill the full surface of a torus.
Thus, surface is an invariant set.

\begin{center}
\includegraphics[width=0.78\textwidth]{lecture3-torus.png}
\end{center}

\newpage
\noindent
When talking about the long-term behavior of trajectories need to be careful
with statements like $\phi_t(\bar x)\to\bar y$ as such limit may not exist,
e.g.\ for periodic orbits.

\medskip
\noindent\underline{Instead use}

\medskip
\noindent\fbox{Definition}

\medskip
\noindent
\underline{$\omega$-limit set} of $\bar x$ is
\[
\omega(\bar x)=\left\{\bar y:\exists\text{ some sequence }t_1,t_2,\ldots
\to\infty\text{ with }\phi_{t_n}(\bar x)\to\bar y\right\}.
\]
\noindent
\underline{$\alpha$-limit set} defined similarly but for $t\to-\infty$.

\begin{center}
\includegraphics[width=0.92\textwidth]{lecture3-limit-set-examples.png}
\end{center}

\noindent
When the orbit $O(\bar x)$ is \underline{bounded}, $\omega(\bar x)$ has
properties:
\[
\begin{array}{ll}
\ast\ \text{non-empty} & \ast\ \text{invariant under }f\\[4pt]
\ast\ \text{closed}    & \ast\ \text{bounded}\\[4pt]
\ast\ \text{connected}&
\end{array}
\]

\noindent\fbox{Example}

\medskip
\noindent
Go back to
\[
\dot r=r(1-r^2),
\qquad
\dot\theta=1.
\]

\begin{center}
\includegraphics[width=0.46\textwidth]{lecture3-limit-cycle-limit-sets.png}
\end{center}

\[
\text{Here}\qquad
\omega(\bar x)=\{r=1\},
\]
\[
\alpha(\bar x)=\{(0,0)\}\text{ for }|\bar x|<1
\qquad\text{or undefined for }|\bar x|>1.
\]

\newpage
\noindent\S 1.4.\ \underline{Topological equivalence of flows}

\bigskip
\noindent
In this course we will sometimes be referring to different dynamical
systems/flows as topologically/geometrically equivalent. Later we'll see that
the flow structure changes at bifurcation.

\bigskip
\noindent\fbox{Definition}\quad\underline{Topological equivalence}

\medskip
\noindent
Two flows $\phi_t^f(\bar x)$ and $\phi_t^g(\bar x)$ are
\underline{topologically equivalent} if there is a homeomorphism
$h(\bar x):E^f\to E^g$ (i.e.\ a continuous bijection with continuous inverse)
and time-increasing function $\tau(\bar x,t)$ such that for any $t,\bar x$
\[
\phi_t^f(\bar x)=h^{-1}\circ\phi_\tau^g\circ h(\bar x)
\]
\[
\text{and}\qquad
\tau(\bar x,t_1+t_2)
=\tau(\bar x,t_1)+\tau\bigl(\phi_{t_1}^f(\bar x),t_2\bigr).
\]

\begin{center}
\includegraphics[width=0.94\textwidth]{lecture3-topological-equivalence.png}
\end{center}

\begin{center}
\textbf{End of lecture 3}
\end{center}

\newpage
\noindent\S 1.4.\ \underline{Topological equivalence of flows}

\bigskip
\noindent
In this course we will sometimes be referring to different dynamical
systems/flows as topologically/geometrically equivalent. Later we'll see that
the flow structure changes at bifurcation.

\bigskip
\noindent\fbox{Definition}\quad\underline{Topological equivalence}

\medskip
\noindent
Two flows $\phi_t^f(\bar x)$ and $\phi_t^g(\bar x)$ are
\underline{topologically equivalent} if there is a homeomorphism
$\bar h(\bar x):E^f\to E^g$ (i.e.\ a continuous bijection with continuous inverse)
and time-increasing function $\tau(\bar x,t)$ such that for any $t,\bar x$
\[
\underline{\phi_t^f(\bar x)=\bar h^{-1}\circ\phi_\tau^g\circ h(\bar x)}
\]
\[
\text{and}\qquad
\underline{\tau(\bar x,t_1+t_2)
=\tau(\bar x,t_1)+\tau\bigl(\phi_{t_1}^f(\bar x),t\bigr)}.
\]

\begin{center}
\includegraphics[width=0.98\textwidth]{lecture4-topological-equivalence.png}
\end{center}

\newpage
\noindent
In words, $\exists$ a `nice' mapping from one phase-space to the other and a map
from one time to another such that the evolutions of the two systems are the same.

\medskip
\noindent
Clearly, $\bar h$ must map \underline{fixed points} into \underline{fixed points}
and \underline{periodic orbits} to \underline{periodic orbits}
(with generally different period).

\medskip
\noindent
If $\tau=t$ we say that the two flows are \underline{topologically conjugate}.

\bigskip
\noindent\fbox{Example}
\[
\begin{aligned}\dot r&=-r\\\dot\theta&=1\end{aligned}
\qquad\text{and}\qquad
\begin{aligned}\dot\rho&=-2\rho\\\dot\psi&=0\end{aligned}
\]
\begin{center}
\includegraphics[height=34mm]{lecture4-spiral.png}
\hspace{35mm}
\includegraphics[height=32mm]{lecture4-radial.png}
\end{center}

\noindent
Flows are topologically equivalent with $\bar h(0,0)=(0,0)$,
$h(r,\theta)=(r^2,\theta+\ln r)$ for $r\neq0$ in polars and
$\tau=t$ (so topologically conjugate).

\medskip
\noindent
To show, solve both: starting at $\bar x=(r_0,\theta_0)$
\[
\underline{\phi_t^f=(r_0e^{-t},\theta_0+t)},\qquad
\underline{\phi_t^g=(\rho_0e^{-2t},\psi_0)}.
\]
\noindent Then
\[
h\circ\phi_t^f
=\bigl(r_0^2e^{-2t},\theta_0+t+\ln(r_0e^{-t})\bigr)
=\bigl(r_0^2e^{-2t},\theta_0+\ln r_0\bigr)
\]
\noindent and
\[
\phi_t^g\circ h=\bigl(r_0^2e^{-2t},\theta_0+\ln r_0\bigr)
\quad\text{-- the same, }\Rightarrow\ \text{Ok}
\]
\[
\text{as }h(\bar x)=(r_0^2,\theta_0+\ln r_0).
\]
\par\noindent\hfill$\blacktriangleleft$

\newpage
\noindent\fbox{Example}
\[
\begin{aligned}\dot r&=0\\\dot\theta&=1\end{aligned}
\qquad\text{and}\qquad
\begin{aligned}\dot r&=0\\\dot\theta&=r+\sin^2\theta\end{aligned}
\]
\noindent
\begin{minipage}[c]{0.25\textwidth}
\includegraphics[width=\linewidth]{lecture4-circular-orbits.png}
\end{minipage}\hfill
\begin{minipage}[c]{0.71\textwidth}
Both are periodic circular orbits, so $\underline{\bar h(\bar x)=\bar x}$,
and time is stretched:
\[
\frac{d\theta}{dt}*(r+\sin^2\theta)=\frac{d\theta}{d\tau},\quad\Rightarrow
\]
\end{minipage}
\[
\underline{\tau=\int_0^t\frac{dt'}{r_0+\sin^2(\theta_0+t')}}
\qquad\begin{gathered}
\text{as first flow is}\\[4pt]
\phi_t(r_0,\theta_0)=(r_0,\theta_0+t).
\end{gathered}
\]
\noindent
Thus, \underline{topologically equivalent} (but not conjugate).
\hfill$\blacktriangleleft$

\bigskip
\noindent\fbox{Definition}

\medskip
\noindent
The vector field $\bar f$ (or system $\dot{\bar x}=\bar f(\bar x)$ or flow
$\phi_t(\bar x)$) is \underline{structurally stable} if $\exists\,\varepsilon>0$
such that $\bar f+\bar\delta$ is topologically equivalent to $f$
$\forall\,\bar\delta(\bar x)$ with
\[
|\bar\delta|+\sum_i\left|\frac{\partial\bar\delta}{\partial x_i}\right|
<\varepsilon.
\]
\noindent
In the first example system is structurally stable, in the second it is not
(a small $\delta\neq0$ will turn $\dot r=0$ into $\dot r=\delta$, which is
no longer a periodic orbit).

\bigskip
\noindent
We will discuss this later in the context of \underline{bifurcations}.

\newpage
\noindent\S 2.\ \underline{Behavior near fixed points}

\medskip
\noindent
Suppose $\bar x_0$ is a fixed point of $\dot{\bar x}=\bar f(\bar x)$,
i.e.\ $\underline{\bar f(\bar x_0)=0}$. If $f$ is smooth near $\bar x_0$, set
\[
\bar y=\bar x-\bar x_0,\quad\Rightarrow\quad
\dot{\bar y}=\overset{\boldsymbol{\cdot}}{(\bar x-\bar x_0)}=\dot{\bar x}
=\bar f(\bar x)=\bar f(\bar y+\bar x_0).
\]
\noindent Taylor expand r.h.s.:
\[
\dot{\bar y}=\underbrace{\bar f(\bar x_0)}_{=0}+\hat A\bar y+O(|y|^2),
\qquad A_{ij}=\frac{\partial f_i}{\partial x_j}
\quad\begin{gathered}\text{-- the Jacobian}\\\text{matrix of }f\text{ at }\bar x_0.
\end{gathered}
\]
\noindent
Thus, near $\bar x_0$ the behavior of the system is approximated by
$\underline{\dot{\bar y}=\hat A\bar y}$.

\bigskip
\noindent\S 2.1\quad\underline{2D flows near fixed points}

\medskip
\noindent
In 2D is a $2\times2$ matrix $\hat A=\begin{pmatrix}a&b\\c&d\end{pmatrix}$
with $\operatorname{Det}(\hat A)=ad-bc$, $\operatorname{Tr}(\hat A)=a+d$.

\medskip
\noindent Eigenvalues satisfy characteristic eq-n
\[
\underline{\lambda^2-\lambda\cdot\operatorname{Tr}(\hat A)+\det(\hat A)=0},
\quad\Rightarrow\quad
\underline{\lambda=\frac12\left[\operatorname{Tr}(\hat A)
\pm\sqrt{(\operatorname{Tr}(\hat A))^2-4\operatorname{Det}(\hat A)}\right]}.
\]
\noindent
By finding eigenvectors (when we can) we can put $\hat A$ into canonical form
\& analyze.

\bigskip
\noindent
\textcircled{1}\quad\underline{Saddle point}: $\underline{\operatorname{Det}(\hat A)<0}$,
$\Rightarrow\lambda_1,\lambda_2$ are \underline{real} and have
\underline{opposite signs}, $\lambda_1\lambda_2<0$.

\medskip
\noindent
\begin{minipage}[c]{0.62\textwidth}
\[
\hat A=\begin{pmatrix}\lambda_1&0\\0&\lambda_2\end{pmatrix},
\quad\text{e.g. }\lambda_1<0<\lambda_2
\]
\[\Updownarrow\]
\[
\dot y_1=\lambda_1y_1,\quad\dot y_2=\lambda_2y_2,
\quad\Rightarrow\quad
\begin{aligned}
y_1&=y_{1,0}e^{\lambda_1t}\\
y_2&=y_{2,0}e^{\lambda_2t}
\end{aligned}
\]
\end{minipage}\hfill
\begin{minipage}[c]{0.35\textwidth}
\includegraphics[width=\linewidth]{lecture4-saddle.png}
\end{minipage}

\begin{center}
\textbf{End of lecture 4}
\end{center}

\newpage
\noindent
Since, eliminating $t$, $y_2\propto y_1^{\lambda_2/\lambda_1}
=\dfrac{1}{y_1^{|\lambda_2/\lambda_1|}}$.
So, orbits are \underline{hyperbolas}!

\bigskip
\noindent
\textcircled{2}\quad\underline{Node}:
$\underline{(\operatorname{Tr}(\hat A))^2>4\operatorname{Det}(\hat A)>0}$,
$\Rightarrow\lambda_1,\lambda_2$ are \underline{real} and either both positive
$(\operatorname{Tr}(\hat A)>0)$ or negative $(\operatorname{Tr}(\hat A)<0)$.

\bigskip
\noindent
\begin{minipage}[t]{0.50\textwidth}
\vspace{0pt}
Same canonical form
\[
\hat A=\begin{pmatrix}\lambda_1&0\\0&\lambda_2\end{pmatrix},
\quad\Rightarrow\quad\underline{y_2\propto y_1^{\lambda_2/\lambda_1}}.
\]
\raggedright
Get different \underline{patterns} depending on
$\underline{\lambda_2/\lambda_1<1\text{ or }>1}$.

\bigskip
\underline{Sign of $\lambda_1,\lambda_2$ sets the}
\underline{direction of the flow}.
\end{minipage}\hfill
\begin{minipage}[t]{0.46\textwidth}
\vspace{0pt}
\includegraphics[width=\linewidth]{lecture5-unstable-nodes.png}
\end{minipage}

\begin{center}
\includegraphics[width=0.78\textwidth]{lecture5-stable-nodes.png}
\end{center}

\newpage
\noindent
\textcircled{3}\quad\underline{Focus (spiral)}:
$\underline{(\operatorname{Tr}(\hat A))^2<4\operatorname{Det}(\hat A)}$
and $\operatorname{det}(\hat A)>0$ so that roots are complex,
$\underline{\lambda=\rho\pm i\omega}$, $\omega\neq0$.

\bigskip
\noindent
\begin{minipage}[c]{0.49\textwidth}
Canonical form
\[
\hat A=\begin{pmatrix}\rho&-\omega\\\omega&\rho\end{pmatrix}.
\]
In polars have
\[
\dot r=\rho r,\quad\dot\theta=\omega,\quad\Rightarrow
\]
\[
\underline{r=r_0e^{\rho t},\quad\theta=\theta_0+\omega t}.
\]
\end{minipage}\hfill
\begin{minipage}[c]{0.49\textwidth}
\includegraphics[width=\linewidth]{lecture5-focus.png}
\end{minipage}

\bigskip
\noindent Other cases are \underline{degenerate}.

\bigskip
\noindent
\textcircled{4}\quad
$\underline{(\operatorname{Tr}(\hat A))^2=4\operatorname{Det}(\hat A)>0}$

\medskip
\noindent
A boundary between \underline{focus \& node}, double eigenvalue
$\underline{\lambda=\frac12\operatorname{Tr}(\hat A)}$. Two cases:

\bigskip
\noindent
\begin{minipage}[c]{0.49\textwidth}
\underline{Stellar node (star)}
\[
\text{if }\hat A=\begin{pmatrix}\lambda&0\\0&\lambda\end{pmatrix}
\text{ is canonical}
\]
\end{minipage}\hfill
\begin{minipage}[c]{0.49\textwidth}
\includegraphics[width=\linewidth]{lecture5-stellar-nodes.png}
\end{minipage}

\newpage
\noindent
\begin{minipage}[c]{0.48\textwidth}
\underline{Improper node}
\[
\text{if }\hat A=\begin{pmatrix}\lambda&1\\0&\lambda\end{pmatrix}
\text{ is canonical}
\]
\centerline{(defective)}
\[
\text{Here}\quad y=y_0e^{\lambda t},
\]
\[
x=x_0e^{\lambda t}+y_0te^{\lambda t},
\qquad\frac{y}{x}=\frac{y_0}{x_0+y_0t}.
\]
\end{minipage}\hfill
\begin{minipage}[c]{0.50\textwidth}
\includegraphics[width=\linewidth]{lecture5-improper-node.png}
\end{minipage}

\bigskip
\noindent
\textcircled{5}\quad
$\underline{\operatorname{Det}(\hat A)>0,\quad\operatorname{Tr}(\hat A)=0}$:
\underline{Centre}

\medskip
\noindent
\begin{minipage}[c]{0.59\textwidth}
\[
\text{Here}\quad\lambda=\pm i\omega,\qquad
\hat A=\begin{pmatrix}0&-\omega\\\omega&0\end{pmatrix}.
\]
\[
\underline{\dot r=0},\quad\underline{\dot\theta=\omega},
\quad\Rightarrow\quad\text{orbits are closed circles}.
\]
\end{minipage}\hfill
\begin{minipage}[c]{0.38\textwidth}
\includegraphics[width=\linewidth]{lecture5-centre.png}
\end{minipage}

\bigskip
\noindent
\textcircled{6}\quad
$\underline{\operatorname{Det}(\hat A)=0,\quad\operatorname{Tr}(\hat A)\neq0}$
-- between \underline{saddle \& node}

\medskip
\noindent
\begin{minipage}[c]{0.48\textwidth}
\[
\lambda_1=0,\quad\lambda_2\neq0,\qquad
\hat A=\begin{pmatrix}0&0\\0&\lambda\end{pmatrix}
\]
\[
y_1=\operatorname{const},\qquad y_2=y_{2,0}e^{\lambda t}
\]
\hfill for $\lambda<0$
\end{minipage}\hfill
\begin{minipage}[c]{0.49\textwidth}
\includegraphics[width=\linewidth]{lecture5-line-fixed-points.png}
\end{minipage}

\newpage
\noindent
\textcircled{7}\quad
$\underline{\operatorname{Det}(\hat A)=0,\quad\operatorname{Tr}(\hat A)=0}$,
$\lambda_{1,2}=0$

\medskip
\noindent
if $\hat A=0$, $\Rightarrow y_1=\operatorname{const},\ y_2=\operatorname{const}$,
$\Rightarrow$ plane of fixed points!

\medskip
\noindent
\begin{minipage}[c]{0.45\textwidth}
if $\hat A=\begin{pmatrix}0&1\\0&0\end{pmatrix}$, this is a case of shear,
\[
\underline{\dot y_1=y_2,\quad\dot y_2=0}.
\]
\end{minipage}\hfill
\begin{minipage}[c]{0.53\textwidth}
\includegraphics[width=\linewidth]{lecture5-shear.png}
\end{minipage}

\medskip
\noindent This zoo can be summarized in a diagram
\begin{center}
\includegraphics[width=0.96\textwidth]{lecture5-classification.png}
\end{center}

\noindent Canonical forms are in special basis
\[
\begin{array}{cll}
(\ast)&\begin{pmatrix}\lambda_1&0\\0&\lambda_2\end{pmatrix}
&\text{-- eigenvectors}\\[12pt]
(\ast)&\begin{pmatrix}\lambda&1\\0&\lambda\end{pmatrix}
&\text{-- generalized eigenvectors}\\[12pt]
(\ast)&\begin{pmatrix}\rho&-\omega\\\omega&\rho\end{pmatrix}
&\text{-- Re \& Im of eigenvectors}
\end{array}
\]

\newpage
\noindent\underline{Notes}

\medskip
\noindent
$(\ast)$ No need to change basis in practice, just work with
$\operatorname{Tr}(\hat A)$, $\operatorname{Det}(\hat A)$ to classify

\medskip
\noindent
$(\ast)$ In original (non-canonical) coordinates shapes can be distorted
\begin{center}
\includegraphics[width=\textwidth]{lecture5-distorted.png}
\end{center}

\noindent
$(\ast)$ Useful to know eigenvectors for saddles/nodes; for focus/centre
cases ratios depend also on $\lambda_{1,2}$, easier to just plot several values

\medskip
\noindent
$(\ast)$ Centres are rare but generic in Hamiltonian systems:
\[
\dot{\bar x}=\begin{pmatrix}\dfrac{\partial H}{\partial y}\\[6pt]
-\dfrac{\partial H}{\partial x}\end{pmatrix}
\quad\text{for some }H(x,y),\quad\bar x=(x,y).
\]
\[
\text{Indeed,}\quad
\hat A=\begin{pmatrix}H_{xy}&H_{yy}\\-H_{xx}&-H_{xy}\end{pmatrix},
\quad\Rightarrow\quad\underline{\operatorname{Tr}(\hat A)=0},\quad\Rightarrow
\]
\[
\text{either}\quad
\underset{\lambda=\pm i\omega}{\underline{\text{centres}}}
\quad\text{or}\quad
\underset{\lambda_1\lambda_2<0}{\underline{\text{saddles}}}.
\]
\[
\text{Also,}\quad
\underline{\frac{dH}{dt}}
=(\dot{\bar x}\cdot\nabla)H
=\dot xH_x+\dot yH_y=H_yH_x-H_xH_y=0.
\]
\noindent
Thus, $H(x,y)=\operatorname{const}$ on trajectories, orbits follow contours
of $H(x,y)$\\
in \underline{autonomous Hamiltonian} systems.

\begin{center}
\textbf{End of lecture 5}
\end{center}

\newpage
% Lecture 6, source PDF page 1 (handwritten page 22).
\noindent\underline{\S 2.2\quad Effect of nonlinear terms}

\medskip
\noindent
How does the behavior of linearization $\dot{\bar y}=\hat A\bar y$ of the
system $\dot{\bar x}=\bar f(\bar x)$ about the fixed point $\bar x_0$
inform us about the behavior of the full nonlinear system around $\bar x_0$?

\medskip
\noindent\fbox{Definition}

\medskip
\noindent
A fixed point $\bar x_0$ is \underline{hyperbolic} if none of the eigenvalues
of the Jacobian of $\bar f(\bar x)$ (matrix $\hat A$) at $\bar x_0$ have zero
real part. It is non-hyperbolic otherwise.

\medskip
\noindent We classify hyperbolic fixed points as
\begin{itemize}
\setlength{\itemsep}{0pt}
\item[$(\ast)$] source if all $\operatorname{Re}(\lambda_i)>0$
\item[$(\ast)$] sink if all $\operatorname{Re}(\lambda_i)<0$
\item[$(\ast)$] saddle otherwise (i.e.\ some $\operatorname{Re}(\lambda_i)>0$, others $<0$).
\end{itemize}
\noindent
We will discuss separately hyperbolic \& non-hyperbolic fixed points.

\medskip
\noindent One can show (e.g.\ see Glendinning, Ch.\ 4) that if
\begin{enumerate}
\setlength{\itemsep}{0pt}
\item[(1)] $\bar x_0$ is hyperbolic and
\item[(2)] the nonlinear corrections of $\bar f$ are $O(|\bar x-\bar x_0|^2)$
\end{enumerate}
\noindent
then the full nonlinear system and the linearized system are topologically conjugate.

\medskip
\noindent\underline{Note}: non-hyperbolic fixed points arise when
\[
\{\operatorname{Det}\hat A\geq0,\ \operatorname{Tr}\hat A=0\},\quad
\{\operatorname{Det}\hat A=0,\ \operatorname{Tr}\hat A\neq0\},
\]
\noindent i.e.\ cases 5, 6, 7 identified earlier.

\newpage
% Lecture 6, source PDF page 2 (handwritten page 23).
\noindent\underline{\S 2.2.1\quad Stable and unstable manifolds}

\medskip
\noindent
For a linearized system we can separate the phase space into different domains
according to different behaviors in time.

\medskip
\noindent\fbox{Definition}

\medskip
\noindent
The \underline{stable}, \underline{unstable} and \underline{centre subspaces}
of the linearization of $\bar f(\bar x)$ near a fixed point $\bar x_0$ are the
three linear subspaces $E^s$, $E^u$ and $E^c$ spanned by subsets of (possibly
generalized) eigenvectors of $\hat A$ whose eigenvalues have real parts
$\operatorname{Re}(\lambda)<0$, $\operatorname{Re}(\lambda)>0$ and
$\operatorname{Re}(\lambda)=0$, respectively.

\medskip
\noindent\underline{Note}: a hyperbolic fixed point has no centre subspace.

\medskip
\noindent\fbox{Example}

\noindent
\begin{minipage}[c]{0.55\textwidth}
\[
\underset{\text{2D}}{\hat A=\begin{pmatrix}-1&0\\0&1\end{pmatrix}},
\quad\Rightarrow\quad
\begin{array}{l}E^s=x\text{-axis}\\E^u=y\text{-axis}\end{array}
\quad\text{1D}
\]
\end{minipage}\hfill
\begin{minipage}[c]{0.43\textwidth}
% Snapshot: Lecture6.pdf, page 2, saddle subspaces.
\includegraphics[width=\linewidth]{lecture6-subspaces.png}
\end{minipage}

\medskip
\noindent\fbox{Theorem}\quad\underline{Stable manifold theorem}

\medskip
\noindent
Suppose $\bar 0$ is a hyperbolic fixed point of $\dot{\bar x}=\bar f(\bar x)$
and $E^s$ and $E^u$ are the stable and unstable subspaces of the linearization
of $\bar f$ near $\bar 0$. Then $\exists$ \underline{stable and unstable manifolds}
$W^s_{\mathrm{loc}}$ \& $W^u_{\mathrm{loc}}$, which have the \underline{same dimension}
as $E^s$, $E^u$ and are \underline{tangent} to $E^s$ \& $E^u$ at $\bar 0$,
such that for $\bar x\neq0$ but in a sufficiently small neighborhood of $\bar 0$,
\[
W^s_{\mathrm{loc}}=\{\bar x:\phi_t(\bar x)\to\bar 0\text{ as }t\to+\infty\},
\]
\[
W^u_{\mathrm{loc}}=\{\bar x:\phi_t(\bar x)\to\bar 0\text{ as }t\to-\infty\}.
\]
\noindent\underline{Proof}: not part of the course, see Glendinning.

\newpage
% Lecture 6, source PDF page 3 (handwritten page 24).
\noindent
The \underline{local} stable manifold $W^s_{\mathrm{loc}}$ can be extended to a
\underline{global invariant manifold} $W^s$ by following the flow backwards in
time starting at $\bar x\in W^s_{\mathrm{loc}}$. Same for $W^u_{\mathrm{loc}}$,
extend to $W^u$ by $t\to+\infty$.

\medskip
\noindent
In $\mathbb R^2$, this theorem guarantees the existence of two trajectories that
approach and leave a saddle fixed point in a tangent fashion.

\medskip
\noindent
\begin{minipage}[c]{0.55\textwidth}
It is often easy in practice to find local approximations to stable and unstable manifolds.

\medskip
Indeed, for a saddle in $\mathbb R^2$ assume that $E^s$ is $x=0$ and $E^u$ is
$y=0$ (can always change the basis).
\end{minipage}\hfill
\begin{minipage}[c]{0.43\textwidth}
% Snapshot: Lecture6.pdf, page 3, local stable and unstable manifolds.
\includegraphics[width=\linewidth]{lecture6-local-manifolds.png}
\end{minipage}

\medskip
\noindent
Then look for $W^s_{\mathrm{loc}}$ in the form $x=S(y)$ with $S(0)=0$ and
$S'(0)=0$. For $W^u_{\mathrm{loc}}$ take $y=U(x)$ with $U(0)=0$, $U'(0)=0$.
These manifolds must satisfy $\dot{\bar x}=\bar f(\bar x)$, i.e.\
$\dot x=\frac{\partial S}{\partial y}\dot y$ and
$\dot y=\frac{\partial U}{\partial x}\dot x$ must solve
$\dot{\bar x}=\bar f(\bar x)$. Thus, for
\[
\begin{pmatrix}\dot x\\\dot y\end{pmatrix}
=\begin{pmatrix}f(x,y)\\g(x,y)\end{pmatrix},
\]
\[
\dot x=\frac{dS}{dy}\dot y=
\underbrace{\frac{dS}{dy}g(S(y),y)=f(S(y),y)}_{\text{(nonlinear) ODE for }S(y)}
\quad\text{and}
\]
\[
\dot y=\frac{dU}{dx}\dot x=
\underbrace{\frac{dU}{dx}f(x,U(x))=g(x,U(x))}_{\text{ODE for }U(x)}.
\]

\newpage
% Lecture 6, source PDF page 4 (handwritten page 25).
\noindent
No need to solve these ODEs, we can simply find a locally convergent series
expansion in the form
\[
S(y)=a_2y^2+a_3y^3+\cdots,\qquad U(x)=b_2x^2+b_3x^3+\cdots
\]
\noindent and solve term by term.

\medskip
\noindent\fbox{Example}
\[
\begin{pmatrix}\dot x\\\dot y\end{pmatrix}
=\begin{pmatrix}x(1-y)\\-y+x^2\end{pmatrix},\qquad
\hat A=\left.\begin{pmatrix}1-y&-x\\2x&-1\end{pmatrix}\right|_{(0,0)}
=\begin{pmatrix}1&0\\0&-1\end{pmatrix},
\]
\[
\det\hat A=-1\ \Rightarrow\ \underline{\text{saddle}},\qquad
\Rightarrow\quad E^s=\{x=0\},\quad E^u=\{y=0\}.
\]
\noindent For $W^u_{\mathrm{loc}}$ take
\[
y=U(x)=a_2x^2+a_3x^3+a_4x^4+\cdots,
\]
\noindent substitute into
\[
\dot y=-y+x^2=\frac{dU}{dx}x(1-y),\quad\Rightarrow
\]
\[
-U(x)+x^2=\frac{dU}{dx}x(1-U(x)),\quad\Rightarrow
\]
\[
\begin{aligned}
x^2-a_2x^2-a_3x^3-a_4x^4+\cdots
&=(2a_2x+3a_3x^2+4a_4x^3+\cdots)\,x\\
&\qquad\times(1-a_2x^2-a_3x^3-a_4x^4+\cdots),\quad\Rightarrow
\end{aligned}
\]
\[
\begin{array}{rclcl}
x^2:&1-a_2=2a_2&&\Rightarrow&a_2=\dfrac13\\[5pt]
x^3:&-a_3=3a_3&&\Rightarrow&a_3=0\\[5pt]
x^4:&-a_4=4a_4-2a_2^2&&\Rightarrow&a_4=\dfrac{2}{45}
\end{array}
\]
\hfill due to $x\mapsto-x$ symmetry!\quad ($a_3=0$)
\[
\text{Thus,}\quad U(x)=\frac{x^2}{3}+\frac{2}{45}x^4+\cdots.
\]

\noindent
\begin{minipage}[c]{0.57\textwidth}
To find $W^s_{\mathrm{loc}}$, $x=S(y)=a_2y^2+a_3y^3+\cdots$,
\[
S(y)(1-y)=\frac{dS}{dy}[-y+S^2(y)],
\]
repeat calculation, find $a_n=0\ \forall n$.
This is caused by symmetry! If $X=-x$ then system the same so
$W^s_{\mathrm{loc}}$ has to be symmetric w.r.t.\ $y$-axis \& tangent to it,
\[
\Rightarrow\quad W^s_{\mathrm{loc}}=y\text{-axis}.
\]
\end{minipage}\hfill
\begin{minipage}[c]{0.41\textwidth}
% Snapshot: Lecture6.pdf, page 4, example phase portrait.
\includegraphics[width=\linewidth]{lecture6-example.png}
\end{minipage}

\begin{center}
\textbf{End of lecture 6}
\end{center}

\newpage
% Lecture 7, source PDF page 1 (handwritten page 26).
\noindent\underline{Note}:
\begin{itemize}
\setlength{\itemsep}{3pt}
\item[$\ast$] If $E^s$ \& $E^u$ are not aligned with the coordinate axes, one can
either transform the coordinates to get alignment, or consider functions $S$ \& $U$
tangent to the relevant eigenvector. E.g.\ for $\bar e=\binom{1}{k}$ use
$y=kx+a_2x^2+a_3x^3+\cdots$.
\item[$\ast$] In higher dimensions use multi-coordinate versions of $S$ \& $U$.
E.g.\ for a 2D $E^s$ in 3D use $S$: $z=S(x,y)$ and
$\dot z=\dot x\dfrac{\partial S}{\partial x}+\dot y\dfrac{\partial S}{\partial y}$
and expand in both coordinates $x,y$.
\end{itemize}

\medskip
\noindent\underline{\S 2.2.2\quad Non-linear terms in non-hyperbolic cases}

\medskip
\noindent
Suppose at least one eigenvalue has $\operatorname{Re}(\lambda)=0$.
Have many possibilities.

\medskip
\noindent (a) $\lambda=\pm i\omega$, linear system has a \underline{centre}.
In nonlinear case multiple possible outcomes, e.g.

\medskip
\noindent
\begin{minipage}[c]{0.70\textwidth}
$\ast\quad\begin{aligned}\dot x&=-y-x^3,\\\dot y&=x-y^3,\end{aligned}$
\quad $\Rightarrow\quad\dot r=-\dfrac{x^4+y^4}{r}\leq0$,
\[
\dot\theta=1-\frac{xy^3-yx^3}{r^2}\to1\quad\text{as }r\to0.
\]
\noindent\hspace*{1em}$H=-\frac12(x^2+y^2)$ for linear.
\end{minipage}\hfill
\begin{minipage}[c]{0.28\textwidth}
% Snapshot citation: Lecture7.pdf, page 1, stable focus.
\includegraphics[width=\linewidth]{lecture7-stable-focus.png}
\end{minipage}

\medskip
\noindent
\begin{minipage}[c]{0.70\textwidth}
$\ast\quad\begin{aligned}\dot x&=-y+x^3,\\\dot y&=x+y^3,\end{aligned}$
\quad $\Rightarrow\quad\dot r=\dfrac{x^4+y^4}{r}\geq0$,
\[
\dot\theta\to1\quad\text{as }r\to0.
\]
\end{minipage}\hfill
\begin{minipage}[c]{0.28\textwidth}
% Snapshot citation: Lecture7.pdf, page 1, unstable focus.
\includegraphics[width=\linewidth]{lecture7-unstable-focus.png}
\end{minipage}

\medskip
\noindent
\begin{minipage}[c]{0.70\textwidth}
$\ast\quad\begin{aligned}\dot x&=-y-2x^2y,\\\dot y&=x+2y^2x\end{aligned}$
\quad Hamiltonian system,
\[
\dot x=\frac{\partial H}{\partial y},\qquad
\dot y=-\frac{\partial H}{\partial x},
\]
with $H=-\frac12(x^2+y^2)-x^2y^2$ --- contours are orbits.
\end{minipage}\hfill
\begin{minipage}[c]{0.28\textwidth}
% Snapshot citation: Lecture7.pdf, page 1, nonlinear centre.
\includegraphics[width=\linewidth]{lecture7-nonlinear-centre.png}
\end{minipage}

\newpage
% Lecture 7, source PDF page 2 (handwritten page 27).
\noindent
(b) $\lambda_1=0$, $\lambda_2\neq0$, \underline{linearized --- line of fixed points}.

\medskip
\noindent Some possibilities:

\medskip
\noindent
\begin{minipage}[c]{0.64\textwidth}
$\ast\quad\begin{aligned}\dot x&=x^2,\\\dot y&=-y,\end{aligned}$
\quad linearized $\begin{pmatrix}0&0\\0&-1\end{pmatrix}$,
\[
\operatorname{Det}\hat A=0,\quad\operatorname{Tr}\hat A=-1.
\]
\end{minipage}\hfill
\begin{minipage}[c]{0.34\textwidth}
% Snapshot citation: Lecture7.pdf, page 2, linearized line of fixed points.
\includegraphics[width=\linewidth]{lecture7-line-fixed-points.png}
\end{minipage}

\noindent
\begin{minipage}[c]{0.57\textwidth}
Full solution:
\[
y\propto\exp\left(\frac1x\right),\quad
\dot x>0,\quad y\propto e^{-t}.
\]
\underline{Note}: $E^c=\{y=0\}$.

\medskip
In the nonlinear case have multiple (non-unique) manifolds tangent to $E^c$
at $(0,0)$ --- no analogue of stable manifold theorem for centre manifolds.
\end{minipage}\hfill
\begin{minipage}[c]{0.41\textwidth}
% Snapshot citation: Lecture7.pdf, page 2, saddle-node.
\includegraphics[width=\linewidth]{lecture7-saddle-node.png}
\end{minipage}

\medskip
\noindent
\begin{minipage}[c]{0.57\textwidth}
$\ast\quad\begin{aligned}\dot x&=x^3,\\\dot y&=-y,\end{aligned}$
\quad $y\propto\exp\left(\dfrac1{2x^2}\right)$.
\end{minipage}\hfill
\begin{minipage}[c]{0.41\textwidth}
% Snapshot citation: Lecture7.pdf, page 2, nonlinear saddle.
\includegraphics[width=\linewidth]{lecture7-nonlinear-saddle.png}
\end{minipage}

\medskip
\noindent
(c) $\lambda_1=\lambda_2=0$, \underline{linearized --- plane or line of fixed points}.

\medskip
\noindent
Almost anything is possible! Change to polars, find rays $\dot\theta=0$,
between them can have

\smallskip
\noindent
% Snapshot citation: Lecture7.pdf, page 2, three sectors in polar coordinates.
\includegraphics[width=\textwidth]{lecture7-polar-sectors.png}

\newpage
% Lecture 7, source PDF page 3 (handwritten page 28).
\noindent\underline{\S 2.3\quad Sketching phase portraits}

\medskip
\noindent\fbox{Example}\quad
$\begin{aligned}\dot x&=x(1-y),\\\dot y&=-y+x^2,\end{aligned}$
\quad $x\mapsto-x$ symmetry, $\Rightarrow$ only $x\geq0$ is enough.

\medskip
\noindent
Fixed points $(0,0)$, $(1,1)$, $(-1,1)$. FP at $(0,0)$ is a \underline{saddle},
studied before. Focus on $(1,1)$.

\medskip
\noindent Linearize:
\[
\hat A=\begin{pmatrix}1-y&-x\\2x&-1\end{pmatrix},\qquad
\hat A_{(1,1)}=\begin{pmatrix}0&-1\\2&-1\end{pmatrix},
\]
\[
\operatorname{Tr}\hat A=-1,\quad\operatorname{Det}\hat A=2,\quad
(\operatorname{Tr}\hat A)^2<4\operatorname{Det}\hat A\quad\Rightarrow\quad
\underline{\text{focus}},
\]
\[
\lambda=\rho\pm i\omega\quad\text{and}\quad
\rho=\tfrac12\operatorname{Tr}\hat A<0\quad\Rightarrow\quad
\underline{\text{stable focus}}.
\]
\noindent
Nullclines (where $\dot x=0$, $\dot y=0$ separately) get signs of $\dot y$, $\dot x$ on them.

\medskip
\noindent
\begin{minipage}[c]{0.46\textwidth}
% Snapshot citation: Lecture7.pdf, page 3, nullclines and flow directions.
\includegraphics[width=\linewidth]{lecture7-nullclines.png}
\end{minipage}\hfill
\begin{minipage}[c]{0.52\textwidth}
% Snapshot citation: Lecture7.pdf, page 3, combined phase portrait.
\includegraphics[width=\linewidth]{lecture7-phase-portrait.png}
\end{minipage}

\medskip
\noindent\underline{General strategy}
\begin{enumerate}
\setlength{\itemsep}{2pt}
\item[(1)] Check for symmetries, find fixed points
\item[(2)] Calculate Jacobian, classify fixed points
\item[(3)] For saddles, find eigenvectors of Jacobian
\item[(4)] Find nullclines (where $\dot x$ or $\dot y=0$) and determine flow direction on them and around
\item[(5)] Use knowledge of stable/unstable manifolds
\item[(6)] Combine these features into a global picture
\item[(7)] \underline{Trajectories cannot cross!}
\item[(8)] Search for periodic orbits (see Chapter 4)
\end{enumerate}

\begin{center}
\textbf{End of lecture 7}
\end{center}

\newpage
% Lecture 8, source PDF page 1 (handwritten page 29).
\noindent\underline{\S 3.\quad Stability}

\medskip
\noindent
If we found a fixed point, or more generally an invariant set $\Lambda$, we know
that starting at $\Lambda$ system will remain there. But what happens if we
slightly perturb the system away from $\Lambda$? Does the system remain near
$\Lambda$? Or does it converge back to or away from $\Lambda$? These are the
issues of system's \underline{stability}, for which there are several definitions.

\medskip
\noindent\underline{\S 3.1\quad Definitions}

\medskip
\noindent
Let $\Lambda$ be an invariant set (a fixed point, periodic orbit, etc.) of a
general (autonomous) system described by a flow $\phi_t$.

\medskip
\noindent\fbox{Definition}\quad\underline{Neighborhood of a set $\Lambda$}

\medskip
\noindent For $\delta>0$ the \underline{neighborhood} of $\Lambda$ is defined as
\[
N_\delta(\Lambda)=\{\bar x:\exists\,\bar y\in\Lambda
\text{ such that }|\bar x-\bar y|<\delta\}.
\]

\noindent\fbox{Definition}\quad\underline{Flow tending to $\Lambda$}

\medskip
\noindent The flow $\phi_t\to\Lambda$ if
\[
\inf_{\bar y\in\Lambda}|\phi_t(\bar x)-\bar y|\to0
\quad\text{as }t\to\infty.
\]

\noindent\fbox{Definition}\quad\underline{Lyapunov stability}

\medskip
\noindent The set $\Lambda$ is \underline{Lyapunov stable} if
$\forall\,\varepsilon>0\ \exists\,\delta>0$ such that if
$\bar x\in N_\delta(\Lambda)$ then
$\phi_t(\bar x)\in N_\varepsilon(\Lambda)\quad\forall\,t\geq0$.

\smallskip
\noindent\textcolor{red}{``Start near, stay near''}

\medskip
\noindent\fbox{Definition}\quad\underline{Quasi-asymptotic stability (QAS)}

\medskip
\noindent The set $\Lambda$ is \underline{quasi-asymptotically stable} if
$\exists\,\delta>0$ such that if $\bar x\in N_\delta(\Lambda)$ then
$\phi_t(\bar x)\to\Lambda$ as $t\to\infty$.

\smallskip
\noindent\textcolor{red}{``Get close eventually''}

\newpage
% Lecture 8, source PDF page 2 (handwritten page 30).
\noindent\fbox{Definition}\quad\underline{Asymptotic stability}

\medskip
\noindent The set $\Lambda$ is \underline{asymptotically stable} if it is both
\underline{Lyapunov stable} and \underline{quasi-asymptotically stable}.

\smallskip
\noindent\textcolor{red}{``Stays near and converges to $\Lambda$''}

\medskip
\noindent\fbox{Example}\quad\underline{$\Lambda$ is LS but not QAS}

\medskip
\noindent
\begin{minipage}[c]{0.62\textwidth}
\[
\begin{pmatrix}\dot x\\\dot y\end{pmatrix}
=\begin{pmatrix}-y\\x\end{pmatrix}
\]
Hamiltonian system with circular orbits. For $\forall\,\varepsilon$ take
$\delta=\zeta\varepsilon$, $0<\zeta<1$, then for $|\bar x|<\delta$ always have
$|\phi_t(\bar x)|<\varepsilon$, $\Rightarrow$ \underline{origin is LS}.
But no convergence to $\bar x=\bar0$, so \underline{not QAS}, and \underline{not AS}.
\end{minipage}\hfill
\begin{minipage}[c]{0.36\textwidth}
% Snapshot citation: Lecture8.pdf, page 2 (handwritten page 30), circular orbits.
\includegraphics[width=\linewidth]{lecture8-circular-orbits.png}
\end{minipage}

\medskip
\noindent\fbox{Example}\quad\underline{$\Lambda$ is QAS but not LS}

\medskip
\noindent
\begin{minipage}[c]{0.42\textwidth}
\[
\begin{pmatrix}\dot r\\\dot\theta\end{pmatrix}
=\begin{pmatrix}r(1-r^2)\\\sin^2\dfrac\theta2\end{pmatrix}
\]
Linearized has one $\lambda=0$, line of fixed points $\parallel$ to $y$-axis
starting at $(1,0)$.

\medskip
Nonlinear terms turn $(1,0)$ into a saddle-node.
\end{minipage}\hfill
\begin{minipage}[c]{0.56\textwidth}
% Snapshot citation: Lecture8.pdf, page 2 (handwritten page 30), homoclinic orbit.
\includegraphics[width=\linewidth]{lecture8-homoclinic-orbit.png}
\end{minipage}

\medskip
\noindent
For $\varepsilon<1$ cannot find $\delta>0$ such that trajectory is in
$N_\varepsilon((1,0))$ --- it will leave it. So, not LS!
But is \underline{QAS} as all trajectories eventually converge at FP $(1,0)$.
So, \underline{not AS}.

\newpage
% Lecture 8, source PDF page 3 (handwritten page 31).
\noindent\fbox{Example}\quad\underline{AS with non-monotonic convergence}

\medskip
\noindent\underline{Stable focus}, elongated in shape.

\medskip
\noindent
\begin{minipage}[c]{0.40\textwidth}
% Snapshot citation: Lecture8.pdf, page 3 (handwritten page 31), elongated stable focus.
\includegraphics[width=\linewidth]{lecture8-elongated-focus.png}
\end{minipage}\hfill
\begin{minipage}[c]{0.58\textwidth}
\underline{LS} since for any $\varepsilon>0$ can always find $\delta>0$ such that
$\phi_t(\bar x)$ with $\bar x\in N_\delta(\bar0)$ always fits in
$N_\varepsilon(\bar0)$. Also \underline{QAS}, obviously, and thus \underline{AS}.
\end{minipage}

\medskip
\noindent But \underline{approaches} FP non-monotonically!

\medskip
\noindent\fbox{Definition}\quad\underline{Basin of attraction}

\medskip
\noindent
If $\Lambda$ is an asymptotically stable invariant set, its \underline{basin of attraction} is
\[
B(\Lambda)=\{\bar x:\phi_t(\bar x)\to\Lambda\text{ as }t\to\infty\}.
\]
If $B(\Lambda)=\mathbb R^n$ then $\Lambda$ is \underline{globally attracting}
(or globally stable).

\medskip
\noindent\fbox{Example}\quad
In last two examples FPs were globally attracting (with some caveats in the saddle-node case).

\medskip
\noindent\fbox{Example}\quad Basin of stability can be \underline{non-trivial}.

\medskip
\noindent
\begin{minipage}[c]{0.53\textwidth}
E.g.\ double potential well with weak friction.

\medskip
\noindent\textcolor{green!55!black}{Hamiltonian contours w/o friction}
\end{minipage}\hfill
\begin{minipage}[c]{0.45\textwidth}
% Snapshot citation: Lecture8.pdf, page 3 (handwritten page 31), double potential well.
\includegraphics[width=\linewidth]{lecture8-double-well.png}
\end{minipage}

\medskip
\noindent For isolated FPs can use linearized system in some cases.

\medskip
\noindent\fbox{Theorem}\quad\underline{Stability of hyperbolic fixed points.}

\medskip
\noindent
If $\bar x_0$ is a hyperbolic \underline{stable focus} or \underline{stable node}
then it is \underline{asymptotically stable} (all hyperbolic sinks are AS).
If $\bar x_0$ is a hyperbolic fixed point with at least one
$\operatorname{Re}(\lambda_i)>0$, then it is \underline{unstable}.

\newpage
% Lecture 8, source PDF page 4 (handwritten page 32).
\noindent\underline{\S 3.2\quad Lyapunov functions}

\medskip
\noindent
We can prove important results about the stability of fixed points of a given
flow $\phi_t(\bar x)$ if we can find some functions that behave in a certain way under the flow.

\medskip
\noindent\fbox{Definition}\quad
A continuously differentiable $(C^1)$ function $V(\bar x):\mathbb R^n\to\mathbb R$
is a \underline{Lyapunov function} for a flow $\phi_t(\bar x)$ on a domain $D$
containing the origin (fixed point) if
\begin{enumerate}
\setlength{\itemsep}{3pt}
\item[(i)] $V(\bar0)=0$ and $V(\bar x)>0$ for $\bar x\neq\bar0$ in $D$ and
\item[(ii)] $\dot V=\bar f\cdot\nabla V\leq0$ in $D$
\quad (or $V(\phi_t(\bar x))\leq V(\bar x)$ for $\forall\,t\geq0$).
\end{enumerate}
\noindent Then we have these 2 theorems.

\medskip
\noindent\fbox{Theorem}\quad\underline{Lyapunov's 1st theorem}

\medskip
\noindent
Suppose that a dynamical system $\dot{\bar x}=\bar f(\bar x)$ has a fixed point at
the origin. If a Lyapunov function exists, then the origin is Lyapunov stable.

\medskip
\noindent\fbox{Proof}\quad
Take $\varepsilon>0$ such that $\{|\bar x|<\varepsilon\}\subseteq D$. Let
\[
\alpha=\min\{V(\bar x):|\bar x|=\varepsilon\}.
\]

\noindent
\begin{minipage}[c]{0.42\textwidth}
% Snapshot citation: Lecture8.pdf, page 4 (handwritten page 32), Lyapunov sublevel set.
\includegraphics[width=\linewidth]{lecture8-lyapunov-sublevel.png}
\end{minipage}\hfill
\begin{minipage}[c]{0.56\textwidth}
Choose $\delta>0$ such that
\[
\{|\bar x|<\delta\}\subset C_{\alpha,\varepsilon}
\]
\[
=\{\bar x:V(\bar x)<\alpha,\ |\bar x|<\varepsilon\},
\]
Then if $|\bar x|<\delta$,
\[
V(\phi_t(\bar x))\leq V(\bar x)<\alpha,
\]
\[
\Rightarrow\quad\underline{|\phi_t(\bar x)|<\varepsilon}.
\]
\end{minipage}

\medskip
\noindent So the origin is \underline{Lyapunov stable (LS)}.

\medskip
\noindent
Intuitively, to leave the ball $|\bar x|<\varepsilon$, a trajectory must first
cross the boundary $|\bar x|=\varepsilon$, where $V\geq\alpha$ by the definition
of $\alpha$. However, $V$ starts below $\alpha$ and is non-increasing along
the trajectory, so it always remains below $\alpha$. The trajectory therefore
cannot reach the boundary and cannot leave the ball.

\begin{center}
\textbf{End of Lecture 8}
\end{center}

\newpage
% Lecture 9 source archive: figures/dynamical-systems/lecture9/Lecture9.pdf.
% Full-page screenshots: lecture9/page-1.png through lecture9/page-5.png.
% Source page 1 (handwritten page 32), saved as lecture9/page-1.png,
% repeats the definition and first theorem already transcribed at the end of Lecture 8.
% Source page 2 (handwritten page 33): lecture9/page-2.png.
\noindent\fbox{Theorem}\quad\underline{Lyapunov's 2nd theorem}

\medskip
\noindent
If, in addition, $\dot V(\bar x)<0$ for $\bar x\neq\bar0$ in $D$, then $V$ is a
\underline{strict Lyapunov function} and the origin is \underline{asymptotically stable}.

\medskip
\noindent\fbox{Proof}\quad
By the first theorem, trajectories starting sufficiently near the origin stay
in a compact neighbourhood $K\subset D$. For such an initial point $\bar x_0$,
$V(\phi_t(\bar x_0))$ decreases and is bounded below by zero, so
\[
V(\phi_t(\bar x_0))\longrightarrow V_0\geq0.
\]
Suppose $V_0>0$. By continuity at the origin, the trajectory is bounded away
from $\bar0$. On the resulting compact subset of $K$, strict negativity and
continuity of $\dot V$ give a constant $\beta>0$ such that
\[
\dot V(\phi_t(\bar x_0))\leq-\beta<0.
\]
But then
\[
V(\phi_t(\bar x_0))
=V(\bar x_0)+\int_0^t\dot V(\phi_s(\bar x_0))\,ds
\leq V(\bar x_0)-\beta t,
\]
which eventually becomes negative, a contradiction. Thus $V_0=0$.
Since $V$ is positive away from $\bar0$ and the trajectory stays in $K$,
$\phi_t(\bar x_0)\to\bar0$. Together with Lyapunov stability, this proves AS.

\medskip
\noindent\fbox{Example}\quad\underline{Stability of the origin}
\[
\begin{pmatrix}\dot x\\\dot y\end{pmatrix}
=\begin{pmatrix}-2x\\x-y\end{pmatrix}.
\]
This is a linear system with Jacobian
\[
A=\begin{pmatrix}-2&0\\1&-1\end{pmatrix},
\qquad \lambda_1=-1,\quad\lambda_2=-2,
\]
so the origin is a \underline{stable node}. Let us illustrate stability using
a Lyapunov function.

\medskip
\noindent Try $V(x,y)=ax^2+by^2$, where $a,b>0$. Clearly $V\geq0$ and $V=0$
only at the origin. Compute
\[
\dot V=V_x\dot x+V_y\dot y
=2ax(-2x)+2by(x-y)
=-2b\left(\frac{2a}{b}x^2-xy+y^2\right).
\]

\newpage
% Source page 3 (handwritten page 34): lecture9/page-3.png.
\noindent
(i) Take $b=8a$. Then
\[
\dot V=-2b\left(\frac{x}{2}-y\right)^2\leq0.
\]
This is a Lyapunov function, but is \underline{not strict}, since $\dot V=0$
on $\{y=x/2\}$, which extends beyond the origin. Still, Lyapunov's first
theorem establishes \underline{Lyapunov stability (LS)}.

\medskip
\noindent
(ii) Take $b=2a$. Then
\[
\dot V=-2b\left[\frac34x^2+\left(\frac{x}{2}-y\right)^2\right]<0
\qquad\text{for }(x,y)\neq(0,0).
\]
A \underline{strict Lyapunov function} has been found, so the second theorem
shows that the origin is \underline{asymptotically stable (AS)}.

\medskip
\noindent Even when the conditions of the second theorem are not fulfilled,
we can sometimes demonstrate asymptotic stability of the origin.

\medskip
\noindent\fbox{Theorem}\quad\underline{LaSalle's Invariance Principle}

\medskip
\noindent
Let $D$ be a compact (closed and bounded), positively invariant set, and let
$V$ be continuously differentiable on a neighbourhood of $D$, with $\dot V\leq0$
in $D$. Then for each $\bar x\in D$ the orbit is bounded, and there exists $c$
such that
\[
\omega(\bar x)\subseteq M_c
=\{\bar z\in D:V(\phi_t(\bar z))=c\ \text{for all }t\geq0\}.
\]
Equivalently, $\phi_t(\bar x)$ tends to an invariant subset of
\[
E=\{\bar z\in D:\dot V(\bar z)=0\}.
\]
% The compact domain in the handwritten statement must be positively invariant;
% this hypothesis also supplies the boundedness used in the proof.

\medskip
\noindent\fbox{Proof}\quad Choose $\bar x\in D$ and let
\[
c=\inf_{t\geq0}V(\phi_t(\bar x)).
\]
Since $D$ is compact and positively invariant, the $\omega$-limit set exists.
Take $\bar y\in\omega(\bar x)$. By definition, there is a sequence $t_n\to\infty$
such that $\phi_{t_n}(\bar x)\to\bar y$. Since $V$ is continuous and
non-increasing along orbits,
\[
V(\bar y)=\lim_{n\to\infty}V(\phi_{t_n}(\bar x))=c.
\]

\newpage
% Source page 4 (handwritten page 35): lecture9/page-4.png.
\noindent
Now show that $V(\phi_t(\bar y))=c$. Suppose instead that
$V(\phi_T(\bar y))<c$ for some $T\geq0$. By continuity of the flow,
\[
V(\phi_{t_n+T}(\bar x))
=V(\phi_T(\phi_{t_n}(\bar x)))
\longrightarrow V(\phi_T(\bar y))<c.
\]
Thus there would be a time $\tau=t_n+T$ with
$V(\phi_\tau(\bar x))<c$, contradicting the definition of $c$.
Hence $V(\phi_t(\bar y))=c$ for all $t\geq0$.
Every point of $\omega(\bar x)$ therefore lies on an orbit on which $V$ is
constant; this gives the invariant subset of $E$ claimed above.

\medskip
\noindent\underline{Notes}
\begin{itemize}
\setlength{\itemsep}{5pt}
\item The $\omega$-limit set need not be a fixed point. This is a more general
statement than Lyapunov's first or second theorem. In practice, we will use
LaSalle's principle to determine the stability of fixed points.
\item If $V$ is positive definite near the origin and the largest invariant
subset of $\{\bar z:\dot V(\bar z)=0\}$ in a suitable compact neighbourhood
is $\{\bar0\}$, then the origin is asymptotically stable. This strengthens
the conclusion of the first theorem without requiring a strict Lyapunov function.
\end{itemize}

\medskip
\noindent\fbox{Example}\quad\underline{Explore stability of the origin}
\[
\begin{pmatrix}\dot x\\\dot y\end{pmatrix}
=\begin{pmatrix}y-xy^2\\-x^3\end{pmatrix}.
\]
Without the $-xy^2$ term, this system would be Hamiltonian, with
\[
H(x,y)=\frac12y^2+\frac14x^4.
\]
Take $V(x,y)=H(x,y)$. Then
\[
\dot V=H_x\dot x+H_y\dot y
=x^3(y-xy^2)+y(-x^3)=-x^4y^2\leq0.
\]
Thus $H$ is a Lyapunov function on $\mathbb R^2$. It is not strict, since
\[
\dot V=0\quad\text{on }\{x=0\}\cup\{y=0\}.
\]

\newpage
% Source page 5 (handwritten page 36): lecture9/page-5.png.
\noindent
\begin{minipage}[c]{0.63\textwidth}
The zero-derivative set extends beyond the origin, so the second theorem
cannot be applied directly. But none of its nonzero points remains in the set:
\[
x=0,\ y\neq0\ \Rightarrow\ \dot x=y\neq0,
\]
\[
y=0,\ x\neq0\ \Rightarrow\ \dot y=-x^3\neq0.
\]
Thus the largest invariant subset of $\{\dot V=0\}$ is $\{\bar0\}$.
By LaSalle's principle, the origin is \underline{asymptotically stable}.
\end{minipage}\hfill
\begin{minipage}[c]{0.34\textwidth}
% Snapshot citation: Lecture9.pdf, page 5 (handwritten page 36),
% cropped from saved full-page screenshot lecture9/page-5.png.
\includegraphics[width=\linewidth]{lecture9/zero-derivative-set.png}
\end{minipage}

\medskip
\noindent To apply the compact-set version of the theorem, take any sublevel
set $\{H\leq k\}$, $k>0$. It is compact and positively invariant since
$\dot H\leq0$.

\medskip
\noindent\underline{Estimating a basin of attraction}

\medskip
\noindent We can use Lyapunov's theorems and LaSalle's principle to establish
a basin of attraction (\underline{domain of stability}, DoS) of a fixed point.
Here is an algorithm:
\begin{enumerate}
\setlength{\itemsep}{9pt}
\item Choose a function $V(\bar x):\mathbb R^n\to\mathbb R$ and a domain $D$
such that
\[
V(\bar x)\geq0\ \text{in }D,\qquad
V(\bar x)=0\ \text{only for }\bar x=\bar0,
\]
and $\dot V(\bar x)\leq0$ in $D$. The latter condition may limit $D$.
\item Find $k>0$ such that the sublevel set
\[
C_k=\{\bar x:V(\bar x)<k\}\subseteq D.
\]
For the compact-set theorem, choose a sublevel region around the origin
whose closure is compact and contained in $D$.
\item Adjust $k$ (or $V$) so that the largest invariant subset of
\[
\overline{C_k}\cap\{\bar z:\dot V(\bar z)=0\}
\]
is the origin. LaSalle's principle then gives
\[
C_k\subseteq\operatorname{DoS}(\bar0).
\]
\item Take different functions $V$ to obtain other such regions $C_k$.
Their union forms a better (larger) estimate of $\operatorname{DoS}(\bar0)$.
\end{enumerate}

\begin{center}
\textbf{End of lecture 9}
\end{center}

\newpage
% Lecture10.pdf, pages 1--5 (handwritten pages 37--41).
\noindent\fbox{Example}\quad\underline{Estimating the domain of stability}
\[
\dot x=x(y^2-1),\qquad \dot y=y(x^2-1).
\]
The fixed points are $(0,0)$ and $(\pm1,\pm1)$. The Jacobian is
\[
A(x,y)=\begin{pmatrix}y^2-1&2xy\\2xy&x^2-1\end{pmatrix}.
\]
At the origin, $A=-I$, so the origin is a stable star node. At each of the
other four fixed points the eigenvalues are $2$ and $-2$, so they are saddles.
The system is symmetric under $(x,y)\mapsto(y,x)$, $x\mapsto-x$, and $y\mapsto-y$.

\medskip
\noindent (1) Try the Lyapunov function
\[
V(x,y)=x^2+a^2y^2,\qquad a>0.
\]
Then
\[
\dot V=2x^2(y^2-1)+2a^2y^2(x^2-1)
=-2(x^2+a^2y^2)+2(1+a^2)x^2y^2.
\]
We need a domain $D$ on which $\dot V<0$ away from the origin. For example,
the region where this holds is described by
\[
\begin{split}
D={}&\left\{(x,y): |x|>\frac{a}{\sqrt{1+a^2}},\quad
 y^2<\frac{x^2}{(1+a^2)x^2-a^2}\right\}\\
&\quad\cup\left\{(x,y): |x|\leq\frac{a}{\sqrt{1+a^2}}\right\}.
\end{split}
\]
Here $\dot V<0$ on $D\setminus\{(0,0)\}$.

\medskip
\noindent (2) Find $k$ so that $C_k=\{(x,y):V(x,y)<k\}\subset D$.
Parametrise a contour of $V$ by
\[
x=\sqrt V\cos\phi,\qquad y=\frac{\sqrt V}{a}\sin\phi.
\]
Substitution gives
\[
\dot V=-2V+\frac{(1+a^2)V^2}{2a^2}\sin^2(2\phi)<0
\quad\text{if}\quad
0<V<\frac{4a^2}{(1+a^2)\sin^2(2\phi)}.
\]
Thus negativity is guaranteed uniformly in $\phi$ whenever
\[
0<V<V_{\max}=\frac{4a^2}{1+a^2}.
\]
\begin{center}
% Original screenshots: Lecture10.pdf, page 1.
\includegraphics[width=.47\textwidth]{lecture10/lyapunov-domain.png}\hfill
\includegraphics[width=.30\textwidth]{lecture10/contour-parametrisation.png}
\end{center}

\newpage
\noindent (3) Hence
\[
C_a=\left\{(x,y):x^2+a^2y^2<\frac{4a^2}{1+a^2}\right\}
\subseteq\operatorname{DoS}(\bar0).
\]
LaSalle's principle is not needed here: $V$ is strict in $C_a$ away from the
origin, so Lyapunov's second theorem applies. More precisely, every point of
$C_a$ lies in a smaller compact sublevel set on which the theorem applies.

\medskip
\noindent (4) Improve the estimate by varying $a$. The ellipses become thin
in the $x$ direction as $a\to0$, are the disc $x^2+y^2<2$ for $a=1$, and become
thin in the $y$ direction as $a\to\infty$:
\[
\frac{x^2}{a^2}+y^2<\frac4{1+a^2},\qquad
x^2+\left(\frac{y}{a^{-1}}\right)^2<\frac4{1+a^{-2}}.
\]
Their union is
\[
\bigcup_{a>0}C_a=\{(x,y):|x|+|y|<2\}.
\]
Indeed, writing $s=a^2>0$, the ellipse inequality becomes
\[
x^2+y^2+\frac{x^2}{s}+s y^2<4,
\]
and $\inf_{s>0}(x^2/s+s y^2)=2|xy|$.
This diamond is the best estimate obtained from this family of Lyapunov
functions. It is not the full domain of stability: on either coordinate
axis, solutions tend to the origin even when $|x|>2$ or $|y|>2$.
\begin{center}
% Original screenshots: Lecture10.pdf, page 2.
\includegraphics[width=.47\textwidth]{lecture10/basin-estimates.png}\hfill
\includegraphics[width=.47\textwidth]{lecture10/basin-phase-portrait.png}
\end{center}

\newpage
\noindent\underline{3.3 Bounding functions}

\medskip
\noindent Another use for functions decreasing along trajectories is to
prove that trajectories eventually enter a bounded region.

\medskip
\noindent\fbox{Definition}\quad\underline{Bounding function}

\medskip
\noindent Suppose $V$ has bounded sublevel sets and there are constants
$M$ and $\delta>0$ such that
\[
\dot V(\bar x)\leq-\delta<0\qquad\text{whenever }V(\bar x)\geq M.
\]
Then all trajectories eventually enter and remain in $\{\bar x:V(\bar x)\leq M\}$.
Such a function is called a bounding function. While $V\geq M$, its value
decreases at least at rate $\delta$, giving finite entry time.

\medskip
\noindent\fbox{Example}\quad Consider the polar system
\[
\dot r=r-r^3(1+k\sin\theta)-r^5,\qquad \dot\theta=1,
\qquad 0<k<1.
\]
Take $V=\frac12r^2$. Then
\[
\dot V=r\dot r=r^2\bigl[1-r^2(1+k\sin\theta)-r^4\bigr]
\leq r^2\bigl[1-r^2(1-k)-r^4\bigr].
\]
For $r^2\geq1/(1-k)$,
\[
\dot V\leq-r^6\leq-\frac1{(1-k)^3}<0.
\]
Thus every trajectory enters the disc
\[
r\leq\frac1{\sqrt{1-k}}
\]
and remains there. What happens within this bounded domain, where an
$\omega$-limit set exists? Is the limit a fixed point or a periodic orbit?
\begin{center}
% Original screenshot: Lecture10.pdf, page 3.
\includegraphics[width=.52\textwidth]{lecture10/bounding-disc.png}
\end{center}

\newpage
\noindent\underline{4. Periodic orbits in $\mathbb R^2$: existence and stability}

\medskip
\noindent We first establish techniques for confirming or excluding the
existence of periodic orbits.

\medskip
\noindent\underline{4.1 Poincar\'e index}

\medskip
\noindent Let
\[
\dot x=f_1(x,y),\qquad \dot y=f_2(x,y),\qquad \bar f=(f_1,f_2).
\]
Away from fixed points, the vector field makes an angle $\psi$ with the
$x$ axis, with $\tan\psi=f_2/f_1$ (using a continuous choice of angle).
Let $\Gamma$ be a simple closed curve containing no fixed point on the curve;
it need not be a trajectory. Traverse $\Gamma$ anticlockwise and measure
the net change in $\psi$. After one circuit the angle changes by an integer
multiple of $2\pi$. This integer is the \underline{Poincar\'e index} $I(\Gamma)$.
\begin{center}
% Original screenshot: Lecture10.pdf, page 4.
\includegraphics[width=.46\textwidth]{lecture10/field-angle.png}
\end{center}
\noindent Its explicit form is
\[
I(\Gamma)=\frac1{2\pi}\oint_\Gamma d\psi
=\frac1{2\pi}\oint_\Gamma
\frac{f_1\,df_2-f_2\,df_1}{f_1^2+f_2^2}.
\]
Any parametrisation of $\Gamma$ may be used to compute this integral. For
example, $df_2=f_{2,x}\,dx+f_{2,y}\,dy$. On a circle,
\[
x=R\cos\phi,\quad y=R\sin\phi,\qquad
 df_2=R(-f_{2,x}\sin\phi+f_{2,y}\cos\phi)\,d\phi.
\]
At a fixed point $f_1=f_2=0$, the angle is undefined; a curve used to compute
$I(\Gamma)$ must avoid such points.

\medskip
\noindent\underline{Properties}
\begin{enumerate}
\setlength{\itemsep}{7pt}
\item $I(\Gamma)$ is unchanged by a continuous deformation of $\Gamma$ as
long as no fixed point is crossed. Indeed, the integral varies continuously
under small deformations, whereas the index is integer-valued, so it must
stay constant.
\item If $\Gamma$ encloses no fixed points, then $I(\Gamma)=0$. Shrink it
to a small contour on which the field is nearly uniform, with zero net
change of angle.
\item The index is additive: if a contour is split into $\Gamma_1$ and
$\Gamma_2$, with the shared joins traversed in opposite directions, then
\[
I(\Gamma)=I(\Gamma_1)+I(\Gamma_2).
\]
The integrals on the shared joins cancel.
\item Reversing time does not change $I(\Gamma)$: replacing $\bar f$ by
$-\bar f$ adds $\pi$ to the angle but does not change its net rotation.
\item For a closed trajectory (a periodic orbit), $I(\Gamma)=+1$, since
the vector field is tangent to the orbit at every point and makes one
positive turn as the curve is traversed anticlockwise.
\item The index of an isolated fixed point is the index of any sufficiently
small simple closed curve encircling that fixed point alone.
\end{enumerate}
\begin{center}
% Original screenshots: Lecture10.pdf, page 5.
\includegraphics[width=.34\textwidth]{lecture10/index-additivity.png}\hfill
\includegraphics[width=.40\textwidth]{lecture10/closed-orbit-index.png}\hfill
\includegraphics[width=.17\textwidth]{lecture10/uniform-field.png}
\end{center}
% Lecture11.pdf page 1 duplicates Lecture10.pdf page 5 (handwritten page 41).
% The unfinished property (7) is continued below, once, in lecture 11.
\begin{center}
\textbf{End of lecture 10}
\end{center}

\newpage
% Lecture11.pdf, pages 2--6 (handwritten pages 42--45 and proof handout).
\noindent\underline{Poincar\'e index: further properties}
\begin{enumerate}
\setcounter{enumi}{6}
\setlength{\itemsep}{8pt}
\item For a curve encircling finitely many isolated fixed points, its index
is the sum of the indices of those fixed points. This follows by deforming
the curve into small loops around the individual points and using additivity.
\item A hyperbolic sink or source has index $+1$, while a hyperbolic saddle
has index $-1$. Non-hyperbolic fixed points must be treated case by case;
for example, a saddle-node has index $0$.
\end{enumerate}
\begin{center}
% Original screenshots: Lecture11.pdf, page 2.
\includegraphics[width=.27\textwidth]{lecture11/index-deformation.png}
\par\medskip
\includegraphics[width=.70\textwidth]{lecture11/node-saddle-indices.png}
\end{center}
\noindent\fbox{Corollary}\quad Any periodic orbit must encircle one or more
fixed points, and the sum of their indices must be $+1$.
This can be used to exclude periodic orbits.

\medskip
\noindent\fbox{Example}\quad\underline{Lotka--Volterra system}
\[
\dot r=r(3-r-2s),\qquad \dot s=s(2-r-s).
\]
The fixed points $(0,0)$, $(3,0)$ and $(0,2)$ are nodes, and $(1,1)$ is a
saddle of index $-1$. A periodic orbit with $r,s>0$ cannot cross the
invariant coordinate axes, so it cannot encircle any of the three nodes.
It can encircle either no fixed point, giving index $0$, or only the saddle,
giving index $-1$. Neither is compatible with the index $+1$ of a periodic
orbit. Thus no such periodic orbit exists.
\begin{center}
\includegraphics[width=.34\textwidth]{lecture11/lotka-volterra-index.png}
\end{center}

\newpage
\noindent\underline{4.2 Dulac's criterion}

\medskip
\noindent\fbox{Theorem}\quad\underline{Dulac's theorem}

\medskip
\noindent If there is a continuously differentiable function $\varphi(x,y)$
such that
\[
\nabla\cdot(\varphi\bar f)>0\qquad\text{everywhere in }D,
\]
or $\nabla\cdot(\varphi\bar f)<0$ everywhere in $D$, where $D\subset\mathbb R^2$
is a simply connected domain, then no periodic orbit lies entirely in $D$.

\medskip
\noindent\fbox{Proof}\quad Suppose $\nabla\cdot(\varphi\bar f)>0$ and a
periodic orbit $\Gamma$ exists. Let $A$ be its interior and $\hat n$ the
outward unit normal. The field is tangent to $\Gamma$, so the divergence
theorem gives
\[
0=\oint_\Gamma\varphi\bar f\cdot\hat n\,d\ell
=\iint_A\nabla\cdot(\varphi\bar f)\,dA>0,
\]
a contradiction. The other sign is treated similarly.
\begin{center}
% Original screenshot: Lecture11.pdf, page 3.
\includegraphics[width=.34\textwidth]{lecture11/dulac-flux.png}
\end{center}
\begin{itemize}
\item Taking $\varphi=1$ gives the \underline{divergence test}: if
$\nabla\cdot\bar f$ is strictly positive or strictly negative throughout a
simply connected domain, there is no periodic orbit in that domain.
\item If $\nabla\cdot(\varphi\bar f)$ has a strict constant sign in a doubly
connected domain, there is at most one periodic orbit, and it must encircle
the hole. Otherwise, applying the divergence theorem to the region between
two periodic orbits gives the same contradiction.
\end{itemize}
\begin{center}
\includegraphics[width=.32\textwidth]{lecture11/doubly-connected-domain.png}
\end{center}

\medskip
\noindent\fbox{Example}\quad\underline{Lotka--Volterra system}
\[
\dot r=r(a-br-cs),\qquad \dot s=s(d-er-fs),\qquad b,f>0.
\]
For $r,s>0$, take $\varphi=1/(rs)$. Then
\[
\nabla\cdot(\varphi\bar f)
=\frac{\partial}{\partial r}\left(\frac{a-br-cs}{s}\right)
+\frac{\partial}{\partial s}\left(\frac{d-er-fs}{r}\right)
=-\frac bs-\frac fr<0.
\]
There are no periodic orbits in the positive quadrant. In the preceding
example, $b=f=1$.

\newpage
\noindent The divergence theorem extends to $\mathbb R^n$ and can exclude
closed invariant hypersurfaces enclosing nonzero volume when the divergence
has a strict constant sign throughout the enclosed region.

\medskip
\noindent\fbox{Example}\quad\underline{Lorenz system}
\[
\dot x=\sigma(y-x),\qquad
\dot y=rx-y-zx,\qquad
\dot z=-bz+xy,\qquad r,\sigma,b>0.
\]
Here
\[
\nabla\cdot\bar f=-\sigma-1-b<0.
\]
Hence there is no closed invariant surface of dimension $2$ enclosing a
nonzero volume. This does not exclude one-dimensional periodic orbits.

\medskip
\noindent\fbox{Theorem}\quad\underline{Gradient criterion}

\medskip
\noindent Suppose there is a positive function $\rho(x,y)$ such that
\[
\rho\bar f=\nabla\psi
\]
for a single-valued potential $\psi$ on a simply connected domain $D$.
Then there are no periodic orbits in $D$.

\medskip
\noindent\fbox{Proof}\quad On a hypothetical nonconstant periodic orbit
$\Gamma$, traversed in the direction of the flow,
\[
0<\oint_\Gamma\rho\bar f\cdot d\bar\ell
=\oint_\Gamma\nabla\psi\cdot d\bar\ell=0.
\]
The last integral vanishes because the potential is single-valued. This
contradiction excludes periodic orbits.

\medskip
\noindent\fbox{Example}\quad Does the system
\[
\dot x=2x+x^2y-y^3,\qquad
\dot y=-2y-xy^2+x^3
\]
have periodic orbits? Notice that
\[
e^{xy}\bar f=\nabla\!\left[e^{xy}(x^2-y^2)\right].
\]
Since $\rho=e^{xy}>0$ and $\psi=e^{xy}(x^2-y^2)$ is single-valued on
$\mathbb R^2$, the gradient criterion excludes periodic orbits.

\newpage
\noindent\underline{4.3 Poincar\'e--Bendixson theorem}

\medskip
\noindent\fbox{Theorem}\quad\underline{Poincar\'e--Bendixson}

\medskip
\noindent If the forward orbit $O^+(\bar x)$ of some point $\bar x$ remains
in a compact set $D\subset\mathbb R^2$ containing no fixed points, then
$\omega(\bar x)$ is a periodic orbit. A sketch proof is given below.
\begin{itemize}
\item This result is specific to $\mathbb R^2$; there is no corresponding
general statement in higher dimensions.
\item Such a trapping region must leave a hole containing a fixed point:
a periodic orbit has index $+1$. In particular, a simply connected domain
containing the whole interior of that orbit cannot be free of fixed points.
\end{itemize}
\begin{center}
% Original screenshot: Lecture11.pdf, page 5.
\includegraphics[width=.21\textwidth]{lecture11/annulus-topology.png}
\end{center}
\noindent\fbox{Example}\quad Consider
\[
\dot x=x-y-2x(x^2+y^2),\qquad
\dot y=y+x-y(x^2+y^2).
\]
Linearising at the origin gives
\[
A=\begin{pmatrix}1&-1\\1&1\end{pmatrix},\qquad
\det A=2,\quad\operatorname{tr}A=2,
\]
so the origin is an unstable focus. Far from the origin the cubic terms
push trajectories inwards, suggesting a periodic orbit.
In polar coordinates,
\[
\dot r=\frac{x\dot x+y\dot y}{r}
=r\bigl[1-r^2(1+\cos^2\theta)\bigr],\qquad
\dot\theta=1+\frac12r^2\sin(2\theta).
\]
The curve $\dot r=0$ away from the origin is
\[
r=\frac1{\sqrt{1+\cos^2\theta}},\qquad 2x^2+y^2=1.
\]
We have $\dot r<0$ for $r>1$ and $\dot r>0$ for $0<r<1/\sqrt2$.
Thus the compact annulus
\[
D=\{(x,y):1/\sqrt2\leq r\leq1\}
\]
is positively invariant. Moreover $\dot\theta\geq1/2>0$ on $D$, so there
are no fixed points in $D$. The Poincar\'e--Bendixson theorem implies the
existence of a periodic orbit there.
\begin{center}
\includegraphics[width=.29\textwidth]{lecture11/trapping-annulus.png}
\end{center}
\noindent\underline{Bonus: uniqueness.}\quad In $D$,
\[
\nabla\cdot\bar f=2-2x^2-5r^2\leq-\frac12-2x^2<0.
\]
By the divergence test for a doubly connected region, there is at most one
periodic orbit in $D$. Hence the periodic orbit just obtained is unique.

\newpage
% Lecture11.pdf page 6: attached printed proof handout.
\noindent\underline{Sketch proof of the Poincar\'e--Bendixson theorem}

\medskip
\noindent The key idea is a \underline{transversal}: a curve segment that
nearby trajectories cross in the same direction. A short segment perpendicular
to the flow can be placed through any non-fixed point.
\begin{center}
% Original screenshot: Lecture11.pdf, page 6.
\includegraphics[width=.57\textwidth]{lecture11/transversal.png}
\end{center}
\noindent\fbox{Lemma}\quad If a nonperiodic trajectory meets a transversal
several times, its successive intersections are arranged monotonically along
the transversal.

\medskip
\noindent After two intersections, the intervening trajectory segment and
the transversal bound a region. Uniqueness of trajectories and the Jordan
curve theorem restrict all further crossings to one direction along the
transversal. This separation into an inside and an outside is where
planarity is essential.
\begin{center}
\includegraphics[width=.95\textwidth]{lecture11/monotone-intersections.png}
\end{center}
\noindent\underline{Proof sketch}
\begin{enumerate}
\setlength{\itemsep}{9pt}
\item Choose $y\in\omega(x)$, then $z\in\omega(y)$. Compactness makes
these $\omega$-limit sets nonempty, and invariance ensures they remain in
$D$, where there are no fixed points. We aim to show that $y$ is periodic.
\item Choose a transversal $\Sigma$ through $z$. Since $z\in\omega(y)$,
the forward orbit of $y$ comes arbitrarily close to $z$. By following these
nearby points forwards or backwards for a small time, obtain two intersections
$y_1,y_2\in O^+(y)\cap\Sigma$ at distinct times.
\item Both $y_1$ and $y_2$ lie in $\omega(x)$ by invariance. The same local
transversal argument gives subsequences of intersections of $O^+(x)$ with
$\Sigma$ tending to $y_1$ and to $y_2$. If $O^+(x)$ is nonperiodic, the lemma
makes its successive intersections monotone, so they cannot have two distinct
limits. Hence $y_1=y_2$, giving a return to the same point at distinct times;
$y$ therefore lies on a periodic orbit. If $O^+(x)$ is already periodic,
the conclusion is immediate.
\end{enumerate}
\begin{center}
\includegraphics[width=.92\textwidth]{lecture11/proof-transversals.png}
\end{center}
\noindent This establishes that $\omega(x)$ contains a periodic orbit.
To obtain the full statement, one also checks that $\omega(x)$ meets any
local transversal in at most one point. Together with connectedness of
$\omega$-limit sets of precompact forward orbits, this shows that $\omega(x)$
is exactly that one periodic orbit. The handout leaves these final details
to the reader.

\begin{center}
\textbf{End of lecture 11}
\end{center}

\newpage
\noindent\underline{4.4 Near-Hamiltonian flows and periodic orbits}

\medskip
\noindent Consider a perturbation of a planar Hamiltonian system:
\[
\begin{pmatrix}\dot x\\\dot y\end{pmatrix}
=\begin{pmatrix}f_1(x,y)+\varepsilon g_1(x,y)\\
f_2(x,y)+\varepsilon g_2(x,y)\end{pmatrix},
\qquad f_1=H_y,\quad f_2=-H_x.
\]
\begin{itemize}
\item For $\varepsilon=0$, $\dot H=0$: trajectories lie on contours of
$H$, and regular closed contours give periodic orbits.
\item For $\varepsilon\ne0$,
\[
\dot H=H_x\dot x+H_y\dot y
=\varepsilon(f_1g_2-f_2g_1).
\]
If a nonconstant periodic orbit $\Gamma$ of period $T$ exists, then
\[
0=\int_0^T\dot H\,dt
=\varepsilon\int_0^T(f_1g_2-f_2g_1)\,dt.
\]
Thus no periodic orbit can lie wholly in a region where
$f_1g_2-f_2g_1$ has a strict constant sign.
\item For $0<\varepsilon\ll1$, $\dot H=O(\varepsilon)$, so trajectories
are weakly perturbed over finite times. The family of Hamiltonian periodic
orbits need not survive. The \underline{energy balance method} tests which
orbits can persist.
\end{itemize}

\noindent Let $\Gamma_0$ be a regular closed orbit with $H=H_0$ and period
$T_0$. Follow a nearby perturbed trajectory $\Gamma_\varepsilon$ for one
return to a local transversal, taking time $T_\varepsilon$. Away from a
separatrix, smooth dependence gives
$\Gamma_\varepsilon=\Gamma_0+O(\varepsilon)$ and
$T_\varepsilon=T_0+O(\varepsilon)$. Hence
\[
\begin{aligned}
\Delta H(\Gamma_\varepsilon)
&=\varepsilon\int_0^{T_\varepsilon}
(f_1g_2-f_2g_1)(\Gamma_\varepsilon(t))\,dt\\
&=\varepsilon M(H_0)+O(\varepsilon^2),\\
M(H_0)&=\int_0^{T_0}(f_1g_2-f_2g_1)(\Gamma_0(t))\,dt.
\end{aligned}
\]
Search for $M(H_0)=0$. A simple zero is robust and gives a nearby periodic
orbit for sufficiently small $\varepsilon$, by the local return map and the
implicit function theorem. A double zero is inconclusive: the neglected
$O(\varepsilon^2)$ terms may remove it.
\begin{center}
\begin{minipage}{.31\textwidth}
\centering
\includegraphics[width=\linewidth]{lecture12/perturbed-orbit.png}\\
{\footnotesize Source: Lecture12.pdf, p.~1 (original crop).}
\end{minipage}\hfill
\begin{minipage}{.51\textwidth}
\centering
\includegraphics[width=\linewidth]{lecture12/energy-roots.png}\\
{\footnotesize Source: Lecture12.pdf, p.~2 (original crop).}
\end{minipage}
\end{center}

\newpage
\noindent\fbox{Example}\quad Consider
\[
\dot x=y,\qquad \dot y=-x+x^2+\varepsilon y(a-x),
\]
or, equivalently,
\[
\ddot x+x-x^2-\varepsilon\dot x(a-x)=0.
\]
For $\varepsilon=0$ the Hamiltonian is
\[
H(x,y)=\frac12(x^2+y^2)-\frac13x^3.
\]
The origin is a centre, $(1,0)$ is a saddle, and the bounded components
of $H=H_0$ for $0<H_0<1/6$ are periodic orbits. The level $H=1/6$
contains a homoclinic orbit to the saddle (a separatrix).
\begin{center}
\includegraphics[width=.48\textwidth]{lecture12/hamiltonian-contours.png}\\
{\footnotesize Source: Lecture12.pdf, p.~2 (original crop).}
\end{center}
\noindent For $\varepsilon\ne0$, $g_1=0$, $g_2=y(a-x)$, so
\[
\dot H=\varepsilon y^2(a-x),\qquad
\Delta H=\varepsilon\int y^2(a-x)\,dt.
\]
The fixed points remain $(0,0)$ and $(1,0)$, with indices $+1$ and $-1$.
Every periodic orbit must enclose the origin and exclude the saddle.
Moreover, $\dot H$ must change sign along it, so it crosses $x=a$.

At extrema of $x$, $\dot x=y=0$ and $\ddot x=x(x-1)$.
At a maximum, $0<x_{\max}<1$; at a minimum, $x_{\min}<0$.
Equality would put the trajectory at a fixed point. Consequently,
\[
x_{\min}<a<x_{\max}<1.
\]
In particular, there are no periodic orbits for $a\geq1$.
\begin{center}
\includegraphics[width=.43\textwidth]{lecture12/orbit-constraints.png}\\
{\footnotesize Source: Lecture12.pdf, p.~3 (original crop).}
\end{center}

\newpage
\noindent\underline{Energy balance for $0<\varepsilon\ll1$}

\medskip
\noindent On the unperturbed orbit $\Gamma_0$, $y\,dt=dx$. Symmetry in
$y$ gives the leading energy change
\[
\begin{aligned}
\Delta H(\Gamma_\varepsilon)&=\varepsilon M(H_0)+O(\varepsilon^2),\\
M(H_0)&=\oint_{\Gamma_0}y(a-x)\,dx\\
&=2\int_{x_{\min}}^{x_{\max}}
\left(2H_0-x^2+\frac23x^3\right)^{1/2}(a-x)\,dx.
\end{aligned}
\]
Here $x_{\min}<0<x_{\max}<1$ are the two turning points of the bounded
orbit, satisfying $H_0=H(x,0)$. The integral is generally evaluated
numerically. The lecture reports no surviving periodic orbit for $a\leq0$
or $a>1/7$, and a surviving periodic orbit for $0<a<1/7$, in this
small-perturbation analysis. The leading energy change increases with $a$,
since
\[
\frac{\partial M}{\partial a}
=2\int_{x_{\min}}^{x_{\max}}
\left(2H_0-x^2+\frac23x^3\right)^{1/2}\,dx>0.
\]
\begin{center}
\includegraphics[width=.67\textwidth]{lecture12/perturbed-phase-portraits.png}\\
{\footnotesize Source: Lecture12.pdf, p.~3 (original crop).}
\end{center}
\begin{center}
\includegraphics[width=.48\textwidth]{lecture12/energy-balance-curves.png}\\
{\footnotesize Source: Lecture12.pdf, p.~3 (original crop).}
\end{center}
\noindent\underline{The leading-order homoclinic threshold.}\quad As
$H_0\uparrow1/6$, the periodic orbit approaches the homoclinic loop, with
$x_{\min}\to-1/2$ and $x_{\max}\to1$. On its upper branch,
\[
\left(\frac13-x^2+\frac23x^3\right)^{1/2}
=\frac1{\sqrt3}(1-x)\sqrt{1+2x}.
\]
Thus the limiting energy balance can be evaluated exactly:
\[
\begin{aligned}
M(1/6)
&=\frac2{\sqrt3}\int_{-1/2}^1(1-x)\sqrt{1+2x}(a-x)\,dx\\
&=\frac65\left(a-\frac17\right).
\end{aligned}
\]
This gives $a=1/7$ as the leading-order threshold. Higher-order terms may
shift the threshold for a fixed nonzero $\varepsilon$; also, the regular
finite-period expansion above is not uniform at the separatrix.

\newpage
\noindent\underline{4.5 Stability of periodic orbits: Floquet analysis}

\medskip
\noindent Consider an autonomous flow $\dot{\bar x}=\bar f(\bar x)$ in
$\mathbb R^n$, with a periodic orbit $\bar x=\bar X(t)$ of period $T$.
Write a nearby solution as
\[
\bar x(t)=\bar X(t)+\bar\eta(t).
\]
Taylor expansion gives
\[
\dot{\bar X}+\dot{\bar\eta}
=\bar f(\bar X+\bar\eta)
=\bar f(\bar X)+A(t)\bar\eta+O(\|\bar\eta\|^2),
\qquad
A_{ij}(t)=\left.\frac{\partial f_i}{\partial x_j}\right|_{\bar X(t)}.
\]
In the linear limit,
\[
\dot{\bar\eta}=A(t)\bar\eta,\qquad A(t+T)=A(t).
\]
This is a linear system with periodic coefficients.

\medskip
\noindent\underline{1. Define the evolution matrix $\Phi(t)$}

\medskip
\noindent By linearity, the map from the initial perturbation to its value
at time $t$ is represented by a matrix:
\[
\Phi(t):\bar\eta(0)\longmapsto\bar\eta(t),\qquad
\boxed{\bar\eta(t)=\Phi(t)\bar\eta(0)}.
\]
In components, following the lecture's notation,
\[
\eta_i(t)=\Phi_{ij}(t)\eta_j(0),\qquad
\Phi_{ij}(0)=\delta_{ij}\quad\Longleftrightarrow\quad\Phi(0)=I.
\]
Repeated indices are summed. The $j$th column of $\Phi(t)$ is the solution
starting from the $j$th unit vector: it records how that initial direction
of perturbation evolves.

Substitute into $\dot\eta_i=A_{ik}(t)\eta_k$:
\[
\dot\Phi_{ij}(t)\eta_j(0)
=A_{ik}(t)\Phi_{kj}(t)\eta_j(0),
\qquad\text{for every }\bar\eta(0).
\]
Thus
\[
\dot\Phi_{ij}=A_{ik}(t)\Phi_{kj},\qquad
\boxed{\dot\Phi(t)=A(t)\Phi(t),\quad\Phi(0)=I}.
\]
This initial-value problem determines $\Phi(t)$; no exponential formula
is needed to define it. The component equations are coupled through the
sum over $k$.

\medskip
\noindent\underline{2. Define the one-period matrix $B$}

\medskip
\noindent Evaluate the evolution matrix after one period:
\[
\boxed{B:=\Phi(T)},\qquad \bar\eta(T)=B\bar\eta(0).
\]
This is the \underline{monodromy matrix}. Unlike the time-dependent
$\Phi(t)$, $B$ is one fixed matrix, describing evolution over a full period.
Periodicity of $A$ and uniqueness imply
\[
\Phi(t+T)=\Phi(t)B.
\]
Indeed, both sides solve the same matrix differential equation and equal
$B$ at $t=0$. Consequently,
\[
\Phi(T)=B,\qquad \Phi(2T)=B^2,\qquad
\boxed{\Phi(mT)=B^m\quad(m=0,1,2,\ldots)}.
\]
Hence $\bar\eta(mT)=B^m\bar\eta(0)$: repeated multiplication by $B$
determines the evolution over successive periods. These conclusions do
not require matrices $A(s)$ and $A(t)$ to commute.

\newpage
\noindent\underline{3. Why does the lecture introduce an integral?}

\medskip
\noindent Having defined $\Phi$ and $B$, the next question is how to
compute them. For a scalar equation,
\[
\dot u=a(t)u
\quad\Longrightarrow\quad
u(t)=\exp\!\left(\int_0^t a(s)\,ds\right)u(0).
\]
The lecture attempts the analogous step for the matrix equation
$\dot\Phi=A(t)\Phi$.

\medskip
\noindent\textcolor{red!70!black}{\fbox{Correction to the lecturer's original formula}}

\medskip
\noindent Lecture12.pdf, p.~4, writes, immediately after ``integrate,
find'',
\[
\Phi_{ij}(t)=B_{ik}(t)\Phi_{kj}(0),\qquad
\widehat B(t)=\exp\!\left(\int_0^t\widehat A(t')\,dt'\right),
\]
with the ordinary matrix exponential defined by
\[
e^{\widehat C}=\sum_{k=0}^{\infty}\frac{\widehat C^{\,k}}{k!}.
\]
Here hats denote matrices. Since $\Phi(0)=I$, the lecture's time-dependent
$\widehat B(t)$ must be $\Phi(t)$, and its constant $\widehat B=\widehat B(T)$
is our monodromy matrix $B$.

\medskip
\noindent\textbf{The issue is the ordinary exponential expression for
$\widehat B(t)$: it does not follow from the hypotheses given.}
In general, matrices at different times need not commute. Periodicity
$A(t+T)=A(t)$ does not imply commutation, and writing the equation in
components does not remove the matrix multiplication or its order.

\medskip
\noindent A sufficient additional condition is
\[
A(s)A(t)=A(t)A(s)\qquad\text{for all }s,t.
\]
Under this condition the lecture's expression is valid, and
\[
\Phi(t)=\exp\!\left(\int_0^t A(s)\,ds\right),\qquad
B=\exp\!\left(\int_0^T A(s)\,ds\right).
\]
For example, if $A$ is constant, $\Phi(t)=e^{tA}$ and $B=e^{TA}$.

\medskip
\noindent In general, the always-valid integral formulation is instead
\[
\boxed{\Phi(t)=I+\int_0^t A(s)\Phi(s)\,ds}.
\]
Equivalently, one can use a time-ordered exponential,
\[
\Phi(t)=\mathcal T\exp\!\left(\int_0^t A(s)\,ds\right),
\]
where $\mathcal T$ preserves the chronological order of the matrix factors.
In practice, solve $\dot\Phi=A(t)\Phi$, $\Phi(0)=I$, up to time $T$ and
then set $B=\Phi(T)$. The original conclusions $\Phi(T)=B$ and
$\Phi(mT)=B^m$ remain valid with this definition of $B$.
% Lecture12.pdf, page 4 ends here; the stability classification continues later.

\begin{center}
\textbf{End of lecture 12}
\end{center}

\end{document}

